\documentclass[12pt,a4paper]{article}

% --- Page geometry as per instructions 
\usepackage[a4paper,left=1in,right=0.5in,top=0.8in,bottom=0.8in]{geometry}

% --- Font setup (Times New Roman look) 
\usepackage{newtxtext,newtxmath}

% --- Line spacing (1.5) 
\usepackage{setspace}
\onehalfspacing

% --- Section heading format (14pt, bold, left-aligned) 
\usepackage{titlesec}
\titleformat{\section}
  {\normalfont\bfseries\fontsize{14}{17}\selectfont}
  {\thesection}
  {1em}
  {}
\titlespacing*{\section}{0pt}{*2}{*1} % Adjust spacing around section titles

% --- Page Border (single line, 3/4 pt width) 
\usepackage{tikz}
\usetikzlibrary{calc}
\usepackage{eso-pic}
\AddToShipoutPictureBG{%
  \begin{tikzpicture}[remember picture,overlay]
    % The border is placed exactly at the margins specified in the geometry package
    \draw[line width=0.75pt]
      ($(current page.north west)+(0.8in,-0.7in)$)
      rectangle
      ($(current page.south east)+(-0.4in,0.7in)$);
  \end{tikzpicture}%
}

% --- Table and Figure packages
\usepackage{graphicx}       % For including images
\usepackage{tabularx}       % For flexible table column widths
\usepackage{caption}        % For table/figure captions
\usepackage{array}          % For table formatting options
\setlength{\extrarowheight}{3pt} % Adds a bit of vertical space to table rows

% --- Other useful packages
\usepackage{ragged2e}       % For \justifying command
\usepackage{parskip}        % To have space between paragraphs instead of indentation
\usepackage[hidelinks]{hyperref} % For clickable links (e.g., in references)

\usepackage{tocloft} % Add this to your preamble
\usepackage{wrapfig}
%\usepackage{tikz} % Already loaded above
\usetikzlibrary{shapes.geometric, shapes.misc}

%\usepackage{algorithm}
\usepackage{algorithmic}
\usepackage{algorithm2e}

\usetikzlibrary{shapes.geometric, shapes.misc, arrows.meta}

% --- Custom column type for left-aligned text in tables
\newcolumntype{L}[1]{>{\RaggedRight\arraybackslash}p{#1}}

% --- Figure, Table, and Equation numbering by section (Chapter-wise)
\usepackage{chngcntr}
\counterwithin{figure}{section}
\counterwithin{table}{section}
\counterwithin{equation}{section}


%-------------------------------------------------------------------------------
%   2. DOCUMENT START
%-------------------------------------------------------------------------------
\begin{document}

%-------------------------------------------------------------------------------
%   TITLE PAGE [cite: 2, 7, 20]
%   No page number on this page
%-------------------------------------------------------------------------------
\begin{titlepage}
    \centering
    \vspace*{0.8cm}

    % --- Project Title (24pt bold as per sample) [cite: 2]
    {\fontsize{22}{28.8}\selectfont\bfseries
    CACHE-AIDED NOMA FOR 6G NETWORKS\\[8pt]}
    
    
    \vspace{1.5cm}

    {\fontsize{14}{16.8}\selectfont\bfseries A PROJECT REPORT} \\[8pt]
    {\fontsize{12}{14.4}\selectfont OF PROJECT-III (PROJ-CS881)}
    
    \vfill % Pushes content to the middle/bottom

    {\fontsize{14}{16.8}\selectfont\bfseries BACHELOR OF TECHNOLOGY}\\[4pt]
    {\fontsize{12}{14.4}\selectfont in}\\[4pt]
    {\fontsize{16}{19.2}\selectfont Information Technology}
    
    {\fontsize{12}{14.4}\selectfont (From \bfseries Maulana Abul Kalam Azad University of Technology,\\West Bengal)}
    
    \vfill

    % --- Submitted by section
    {\fontsize{12}{14.4}\selectfont SUBMITTED BY:}\\[12pt]
    {\fontsize{14}{16.8}\selectfont % 14pt for names
    \begin{tabular}{c}
        Akash Mandal (13000222028) \\
        Arpita Guchhait (13000222031) \\
        Manojit Pal (13000222036) \\
        Vaibhav Anand (13000222037) \\
    \end{tabular}}
    
    \vspace{1.5cm}
    
    \vfill
    
    \begin{figure}[htbp]
    {\centering {\includegraphics[scale=0.22]{TMSL.jpg}}\par}
    \end{figure}

    % --- Department and Address
    {\large\fontsize{14}{16.8}\selectfont\bfseries
    Department of Information Technology\\
    Techno Main Salt Lake\\
    EM-4/1, Salt Lake City, Kolkata - 700091\\[2pt]
    }
    
    \vspace{0.8cm}
\end{titlepage}

%-------------------------------------------------------------------------------
%   CERTIFICATE SECTION (Page ii)
%   Font: Times New Roman, Size 14, Double line spacing
%-------------------------------------------------------------------------------
\clearpage
\pagenumbering{roman}
\setcounter{page}{2}

\begin{center}
    {\fontsize{16}{19.2}\selectfont\bfseries BONAFIDE CERTIFICATE}\\[1cm]
\end{center}

\begin{doublespace}
{\fontsize{14}{16.8}\selectfont
This is to certify that the project report entitled \textbf{``CACHE-AIDED NOMA FOR 6G NETWORKS''} submitted by \textbf{Akash Mandal (13000222028), Arpita Guchhait (13000222031), Manojit Pal (13000222036), and Vaibhav Anand (13000222037)} in partial fulfillment of the requirements for the award of the degree of Bachelor of Technology in Information Technology, Maulana Abul Kalam Azad University of Technology, West Bengal, is a record of bonafide work carried out by them under our supervision and guidance. The results embodied in this report have not been submitted to any other University or Institute for the award of any degree or diploma.
}
\end{doublespace}

\vspace{2cm}

\noindent
\begin{minipage}[t]{0.45\textwidth}
    \begin{center}
        \rule{6cm}{1pt}\\
        \textbf{Signature of the Supervisor(s)}\\
        Dr. Biswajit Ghosh \& Mr. Abdul Aziz\\
        Department of Information Technology\\
        Techno Main Salt Lake
    \end{center}
\end{minipage}
\hfill
\begin{minipage}[t]{0.45\textwidth}
    \begin{center}
        \rule{6cm}{1pt}\\
        \textbf{Signature of the HOD}\\
        Department of Information Technology\\
        Techno Main Salt Lake
    \end{center}
\end{minipage}

\vspace{1.5cm}
\noindent
Date: ........................

\thispagestyle{plain}


%-------------------------------------------------------------------------------
%   ACKNOWLEDGMENT SECTION (Page iii)
%-------------------------------------------------------------------------------
\clearpage

\begin{center}
    {\fontsize{16}{19.2}\selectfont\bfseries ACKNOWLEDGMENT}\\[1cm]
\end{center}

\begin{onehalfspacing}
{\fontsize{12}{14.4}\selectfont
We would like to express our deep gratitude to our supervisors, \textbf{Dr. Biswajit Ghosh} and \textbf{Mr. Abdul Aziz}, for their patient guidance, enthusiastic encouragement, and useful critiques of this research work. Their willingness to give their time so generously has been very much appreciated.

We would also like to thank the Head of the Department, Department of Information Technology, Techno Main Salt Lake, for providing us with the necessary facilities and support to complete this project.

Finally, we wish to thank our parents and friends for their continuous support and encouragement throughout our study and project work.
}
\end{onehalfspacing}

\vspace{1.5cm}
\begin{flushright}
\textbf{Akash Mandal} (13000222028)\\
\textbf{Arpita Guchhait} (13000222031)\\
\textbf{Manojit Pal} (13000222036)\\
\textbf{Vaibhav Anand} (13000222037)
\end{flushright}

\thispagestyle{plain}

%-------------------------------------------------------------------------------
%   ABSTRACT SECTION (Page iv)
%   Font: Times New Roman, Size 12, Single line spacing
%   Page numbering: Roman numerals (typically page iv)
%-------------------------------------------------------------------------------

\clearpage
\section*{Abstract}
\addcontentsline{toc}{section}{Abstract}

\begin{singlespace}
\noindent
This project proposes a Deep Reinforcement Learning (DRL) based proactive caching framework for Cache-aided Non-Orthogonal Multiple Access (NOMA) downlink systems targeting 6G networks. While caching and NOMA are promising complementary technologies, traditional reactive caching policies fail to optimize for Cache-aided Interference Cancellation (CIC) opportunities. The proposed scheme employs a Double Dueling Deep Q-Network (DQN) agent with a novel binary action space to make online cache placement decisions. By incorporating a CIC-aware reward function, the agent is explicitly incentivized to proactively cache files that enable interference suppression for paired NOMA users. Simulated under a realistic mixed fading channel model with imperfect Successive Interference Cancellation (SIC), results over 5 million training steps and 100 Monte Carlo evaluation runs demonstrate exceptional performance. The proposed DQN scheme achieves a cache hit rate of 70.38\%, a CIC benefit rate of 3.78 (a 5.28x improvement over the best dynamic baseline), and a sum-rate of 18.17 bps/Hz at an SNR of 30 dB. Furthermore, the proposed policy reduces system outage probability by 14x compared to the no-caching baseline, performing within 0.54\% of the oracle-optimal Top-K policy without relying on prior knowledge of the content popularity distribution.

\vspace{0.5cm}

\noindent
\textbf{Keywords:} Cache-Aided NOMA, Deep Reinforcement Learning, 6G Networks, Proactive Caching, Cache-aided Interference Cancellation (CIC), Double Dueling DQN, Binary Action Space.
\end{singlespace}

\thispagestyle{plain}

%-------------------------------------------------------------------------------
%   TABLE OF CONTENTS
%   Font: Times New Roman, Size 12, One-and-a-half spacing
%   Page numbering: Roman numerals (continued from abstract)
%-------------------------------------------------------------------------------

\clearpage
% Page numbering continues in roman numerals automatically

% Customize ToC appearance
\renewcommand{\cfttoctitlefont}{\normalfont\bfseries\fontsize{14}{17}\selectfont}
\renewcommand{\cftaftertoctitle}{\hfill}
\renewcommand{\contentsname}{TABLE OF CONTENTS} % All caps to match IoT format

% Adjust spacing for entries
\setlength{\cftsecnumwidth}{2.5em}
\setlength{\cftsecindent}{0pt}
\setlength{\cftbeforesecskip}{6pt}
\setlength{\cftsubsecnumwidth}{3.5em}
\setlength{\cftsubsecindent}{2.5em}

% Add dots for sections (optional, matches many report formats)
\renewcommand{\cftsecleader}{\cftdotfill{\cftdotsep}}

\tableofcontents
\thispagestyle{plain}

%-------------------------------------------------------------------------------
%   LIST OF TABLES
%   Font: Times New Roman, Size 12, One-and-a-half spacing
%   Captions must match exactly as they appear in text
%   Page numbering: Roman numerals (continued)
%-------------------------------------------------------------------------------

\clearpage
% Page numbering continues automatically in roman numerals

% Customize List of Tables title
\renewcommand{\cfttabpresnum}{Table }
\renewcommand{\cfttabnumwidth}{5em}
\renewcommand{\listtablename}{LIST OF TABLES} % All caps

% Add to table of contents
\addcontentsline{toc}{section}{List of Tables}

\listoftables
\thispagestyle{plain}

%-------------------------------------------------------------------------------
%   LIST OF FIGURES
%   Font: Times New Roman, Size 12, One-and-a-half spacing
%   Captions must match exactly as they appear in text
%   Page numbering: Roman numerals (continued)
%-------------------------------------------------------------------------------

\clearpage
% Page numbering continues automatically in roman numerals

% Customize List of Figures title
\renewcommand{\cftfigpresnum}{Figure }
\renewcommand{\cftfignumwidth}{5.5em}
\renewcommand{\listfigurename}{LIST OF FIGURES} % All caps

% Add to table of contents
\addcontentsline{toc}{section}{List of Figures}

\listoffigures
\thispagestyle{plain}

%-------------------------------------------------------------------------------
%   LIST OF SYMBOLS, ABBREVIATIONS AND NOMENCLATURE
%   Font: Times New Roman, Size 12, One-and-a-half spacing
%-------------------------------------------------------------------------------

\clearpage

\section*{List of Symbols, Abbreviations and Nomenclature}
\addcontentsline{toc}{section}{List of Symbols, Abbreviations and Nomenclature}

\subsection*{Abbreviations}

\begin{tabular}{@{}ll@{}}
\textbf{5G} & Fifth Generation \\
\textbf{6G} & Sixth Generation \\
\textbf{AAT} & Average Access Time \\
\textbf{BER} & Bit Error Rate \\
\textbf{BS} & Base Station \\
\textbf{CSI} & Channel State Information \\
\textbf{D2D} & Device-to-Device \\
\textbf{DQN} & Deep Q-Network \\
\textbf{DRL} & Deep Reinforcement Learning \\
\textbf{eMBB} & Enhanced Mobile Broadband \\
\textbf{LFU} & Least Frequently Used \\
\textbf{LRU} & Least Recently Used \\
\textbf{MEC} & Mobile Edge Computing \\
\textbf{ML} & Machine Learning \\
\textbf{mMTC} & Massive Machine-Type Communication \\
\textbf{NOMA} & Non-Orthogonal Multiple Access \\
\textbf{OMA} & Orthogonal Multiple Access \\
\textbf{QoS} & Quality of Service \\
\textbf{SIC} & Successive Interference Cancellation \\
\textbf{SINR} & Signal-to-Interference-plus-Noise Ratio \\
\textbf{SNR} & Signal-to-Noise Ratio \\
\textbf{URLLC} & Ultra-Reliable Low-Latency Communication \\
\end{tabular}

\subsection*{Symbols}

\begin{tabular}{@{}ll@{}}
$\alpha$ & Zipf skewness parameter \\
$\epsilon$ & Exploration rate in epsilon-greedy policy \\
$\tau$ & Soft update parameter for target network \\
$\gamma$ & Signal-to-Interference-plus-Noise Ratio (SINR) \\
$P_t$ & Transmission power \\
$P_1, P_2$ & Power allocation for user 1 and user 2 \\
$h$ & Channel coefficient \\
$N_0$ & Noise power spectral density \\
$R$ & Data rate (bps/Hz) \\
$C$ & Cache capacity \\
$N$ & Number of users \\
$M$ & Content catalog size \\
$Q(s,a)$ & Q-value function \\
$s$ & State representation \\
$a$ & Action \\
$r$ & Reward signal \\
\end{tabular}

\thispagestyle{plain}

%-------------------------------------------------------------------------------
%   CHAPTER-I DIVIDER PAGE
%   Matches IOT.pdf page 8 format
%   Page numbering starts with Arabic numerals (page 1)
%-------------------------------------------------------------------------------

\clearpage
\pagenumbering{arabic}
\setcounter{page}{1}

\begin{center}
    \vspace*{10cm}
    
    {\fontsize{16}{19.2}\selectfont\bfseries
    Chapter -- I}
    
    \vspace{1cm}
    
    {\fontsize{16}{19.2}\selectfont\bfseries
    INTRODUCTION}
    
    \vspace*{\fill}
\end{center}

\thispagestyle{plain}

%-------------------------------------------------------------------------------
%   CHAPTER 1: INTRODUCTION CONTENT (Compact - 4 Pages)
%   Following IOT.pdf structure with essential formulas
%-------------------------------------------------------------------------------

\clearpage

\section{Introduction}

\subsection{Introduction to 6G Networks}

Wireless communication has evolved through multiple generations: 1G (analog voice), 2G (digital cellular), 3G (mobile internet), 4G/LTE (mobile broadband up to 100 Mbps), and 5G supporting Enhanced Mobile Broadband (eMBB, 20 Gbps), Ultra-Reliable Low-Latency Communication (URLLC, <1 ms latency), and Massive Machine-Type Communication (mMTC, 1 million devices/km²).

Emerging applications like holographic communications, tactile internet, and brain-computer interfaces demand capabilities beyond 5G, driving research toward Sixth Generation (6G) networks by 2030. 6G targets 1 Terabit per second data rates (100× faster), 0.1-1 ms latency, 10 million devices/km² density, 100× energy efficiency, and 5-10× spectrum efficiency improvement.

Critical challenges include spectrum scarcity as traffic reaches 5,000 exabytes/month by 2030, backhaul congestion where 20\% of content accounts for 80\% of requests, inability to meet ultra-low latency with cloud architectures, energy sustainability requirements, and user fairness for cell-edge users.

\subsection{Introduction to Non-Orthogonal Multiple Access (NOMA)}

Traditional schemes (TDMA, FDMA, CDMA, OFDMA) allocate dedicated orthogonal resources to each user, causing inefficient spectrum utilization. Non-Orthogonal Multiple Access (NOMA) enables multiple users to share identical time-frequency resources via power-domain multiplexing and Successive Interference Cancellation (SIC).

\subsubsection{NOMA Principles}

NOMA allocates higher power to far users (weak channels) and lower power to near users (strong channels). For two users, the transmitted signal is:
\begin{equation}
x(t) = \sqrt{P_1} \cdot s_1(t) + \sqrt{P_2} \cdot s_2(t)
\end{equation}
where $P_1 > P_2$ for far/near users, satisfying $P_1 + P_2 = P_{\text{total}}$. Common allocation: $P_1 = \alpha \cdot P_{\text{total}}$, $P_2 = (1-\alpha) \cdot P_{\text{total}}$ with $\alpha \in [0.7, 0.8]$.

The received signal is: $y_i(t) = h_i \cdot x(t) + n_i(t)$ where $h_i$ is channel coefficient and $n_i(t)$ is noise.

\subsubsection{Successive Interference Cancellation}

Near user (user 2) performs SIC: first decodes far user's signal with SINR:
\begin{equation}
\gamma_{2 \rightarrow 1} = \frac{P_1 |h_2|^2}{P_2 |h_2|^2 + N_0}
\end{equation}

After subtraction with imperfection factor $\eta \in [0.05, 0.15]$, decodes own signal:
\begin{equation}
\gamma_2 = \frac{P_2 |h_2|^2}{\eta P_1 |h_2|^2 + N_0}
\end{equation}

Far user decodes directly: $\gamma_1 = \frac{P_1 |h_1|^2}{P_2 |h_1|^2 + N_0}$

Achievable rate: $R_i = B \log_2(1 + \gamma_i)$ where $B$ is bandwidth.

NOMA achieves 30-50\% spectral efficiency gain and 2-3× user capacity but faces challenges: imperfect SIC ($\eta = 0.05-0.15$), CSI errors (5-10\%), 2-3× computational complexity, error propagation, and complex multi-user optimization.

\subsection{Introduction to Mobile Edge Caching}

Mobile Edge Caching (MEC) stores popular content at base stations, reducing latency from 50-200 ms (backhaul) to 1-5 ms (local cache).

Content requests follow Zipf distribution:
\begin{equation}
P(k) = \frac{k^{-\alpha}}{\sum_{i=1}^{N} i^{-\alpha}}
\end{equation}
where $k$ is rank, $N$ is library size, and $\alpha \in [0.6, 1.2]$ is skewness. With $\alpha = 1.0$, top 10\% content receives 60\% requests.

Traditional policies: \textbf{LRU} (evicts least recently used), \textbf{LFU} (evicts least frequently used), \textbf{Static Top-K} (pre-caches top $K$ items). These reactive policies cannot anticipate future popularity.

Performance metrics: Cache Hit Ratio $= \frac{\text{Hits}}{\text{Total Requests}}$ and Average Access Time:
\begin{equation}
\text{AAT} = p_{\text{hit}} \cdot T_{\text{cache}} + (1 - p_{\text{hit}}) \cdot T_{\text{backhaul}}
\end{equation}

\subsection{Introduction to Deep Reinforcement Learning}

Deep Reinforcement Learning (DRL) combines RL decision-making with deep neural networks for learning optimal policies from high-dimensional observations.

RL is a Markov Decision Process (MDP) with states $\mathcal{S}$, actions $\mathcal{A}$, transitions $\mathcal{P}$, rewards $\mathcal{R}$, and discount $\gamma$. Agent maximizes cumulative reward:
\begin{equation}
G_t = \sum_{k=0}^{\infty} \gamma^k r_{t+k}
\end{equation}

Action-value function: $Q^\pi(s,a) = \mathbb{E}_\pi[ G_t \mid s_t = s, a_t = a ]$

\subsubsection{Deep Q-Networks (DQN)}

DQN approximates $Q(s,a)$ using neural networks $Q(s,a;\theta)$ with key innovations:

\textbf{Experience Replay:} Stores transitions $(s_t, a_t, r_t, s_{t+1})$ in buffer $\mathcal{D}$, trains on random mini-batches.

\textbf{Target Network:} Separate network $Q(s,a;\theta^-)$ with soft updates: $\theta^- \leftarrow \tau \theta + (1-\tau) \theta^-$ where $\tau = 0.001-0.01$.

\textbf{Epsilon-Greedy:} Random action with probability $\epsilon$, greedy with $1-\epsilon$. $\epsilon$ anneals from 1.0 to 0.01.

\textbf{Loss Function:}
\begin{equation}
L(\theta) = \mathbb{E}_{(s,a,r,s') \sim \mathcal{D}} \left[ \left( r + \gamma \max_{a'} Q(s',a';\theta^-) - Q(s,a;\theta) \right)^2 \right]
\end{equation}

\textbf{Dueling DQN:} Decomposes Q-values:
\begin{equation}
Q(s,a;\theta) = V(s;\theta_V) + A(s,a;\theta_A) - \frac{1}{|\mathcal{A}|}\sum_{a'} A(s,a';\theta_A)
\end{equation}

\textbf{Prioritized Replay:} Samples transitions by TD error: $p_i = \frac{|\delta_i|^\alpha}{\sum_j |\delta_j|^\alpha}$

For caching, DRL maps states (cache contents, requests), actions (cache decisions), and rewards (hits, latency) to learn optimal policies adapting to temporal dynamics.

\subsection{Motivation of the Work}

Traditional orthogonal schemes face spectral constraints as density increases. Reactive caching policies cannot anticipate popularity shifts. Separate optimization misses synergies between caching and NOMA. Static power allocation doesn't adapt to cache status.

Cache-aided NOMA integration creates complementary benefits: caching reduces latency and offloads 40-60\% backhaul traffic, while NOMA enhances spectral efficiency by 30-50\% and increases capacity 2-3 times. Together, they enable efficient simultaneous delivery of cached content via NOMA. Deep reinforcement learning provides the intelligence layer for optimal joint decisions.

\subsubsection{Research Objectives}

This project develops an intelligent cache-aided NOMA framework using DQN with specific objectives:

\textbf{Objective 1:} Design a novel NOMA-aware state representation (14-dimensional) for the DRL agent that jointly encodes content popularity, cache occupancy, continuous channel statistics, and recent NOMA pairing history.

\textbf{Objective 2:} Present a binary action space redesign that reduces the traditional DQN eviction action space from $\mathcal{O}(N_{\text{files}})$ to $\mathcal{O}(1)$, making DRL remarkably tractable for large, real-world content catalogs.

\textbf{Objective 3:} Propose a specialized CIC-incentivized immediate reward function that explicitly rewards cache placement decisions enabling interference cancellation for paired NOMA users.

\textbf{Objective 4:} Implement a comprehensive NOMA simulator with realistic mixed fading channels, path loss, configurable power allocation, and SIC with practical imperfection factors ($\eta = 0.05-0.15$).

\textbf{Objective 5:} Conduct an extensive simulation study comparing the DQN scheme against four baselines (Top-K Zipf, LRU, LFU, Random) and a No-Cache reference over a full signal-to-noise ratio (SNR) sweep to evaluate cache hit ratio, outage probability, and system sum-rate.

By achieving these objectives, this project establishes AI-driven cache-aided NOMA, optimized for Cache-aided Interference Cancellation (CIC), as a highly viable and scalable strategy for meeting 6G network demands.

\thispagestyle{plain}

%-------------------------------------------------------------------------------
%   CHAPTER-II DIVIDER PAGE
%   Matches IOT.pdf page 13 format
%   Page numbering continues from Chapter 1
%-------------------------------------------------------------------------------

\clearpage

\begin{center}
    \vspace*{10cm}
    
    {\fontsize{16}{19.2}\selectfont\bfseries
    Chapter -- II}
    
    \vspace{1cm}
    
    {\fontsize{16}{19.2}\selectfont\bfseries
    LITERATURE REVIEW}
    
    \vspace*{\fill}
\end{center}

\thispagestyle{plain}

%-------------------------------------------------------------------------------
%   CHAPTER 2: LITERATURE REVIEW CONTENT (Condensed - 2 Pages)
%   Following IOT.pdf structure - paragraph format without subsections
%   Proper citations using refs.bib
%-------------------------------------------------------------------------------

\clearpage

\section{Literature Review}

The research landscape surrounding cache-aided wireless networks and non-orthogonal multiple access has evolved significantly over the past decade. Early empirical investigations established foundational principles for modern caching strategies. Breslau et al. \cite{b_breslau_zipf} demonstrated that web content popularity follows Zipf-like distributions, where request probability for the $k$-th most popular item is inversely proportional to its rank raised to a power. This observation became a cornerstone for cache placement analysis, with typical skewness parameters ranging from 0.6 to 1.2, implying that a small cache storing popular content can satisfy a disproportionately large fraction of requests.

Building on these insights, researchers explored architectural innovations to push content closer to end users. Golrezaei et al. \cite{b_golrezaei} introduced femtocaching, where small-cell base stations with local storage proactively cache popular content, demonstrating substantial backhaul offloading and latency reduction. Ji et al. \cite{b_ji_d2d} investigated caching limits in Device-to-Device wireless networks, deriving information-theoretic bounds on cache-aided throughput. A major theoretical breakthrough emerged from Maddah-Ali and Niesen \cite{b_maddah_ali}, who introduced coded caching revealing that coordinated placement with coded multicast delivery achieves multiplicative throughput gains scaling with the number of users, demonstrating that caching benefits extend beyond simple hit-rate improvements to fundamental traffic reductions.

In parallel, Non-Orthogonal Multiple Access emerged as a promising physical layer technique for 5G and beyond. Dai et al. \cite{b_dai_noma} provided comprehensive treatment of power-domain NOMA, detailing power allocation strategies, Successive Interference Cancellation procedures, and performance under Rayleigh fading. Their work highlighted NOMA's potential to improve spectral efficiency by 30-50\% while supporting massive connectivity. Ding et al. \cite{b_ding_survey} presented an extensive survey comparing NOMA with Orthogonal Frequency Division Multiple Access, identifying practical challenges including imperfect SIC with residual interference typically 5-15\% of cancelled signal power, channel estimation errors, and increased receiver complexity requiring iterative decoding.

The convergence of edge caching and NOMA has attracted considerable attention for their complementary benefits in addressing spectrum scarcity and latency simultaneously. Bepari et al. \cite{b_survey_bepari} conducted a comprehensive survey on cache-aided NOMA for sixth-generation networks, reviewing architectures, performance analysis, optimization frameworks, and applications including UAV communications and mobile edge computing. They identified cache hit ratio, average access time, energy efficiency, and outage probability as key performance indicators. Zhang et al. \cite{b_zhang_stochastic} developed stochastic geometry frameworks modeling cache-enabled NOMA with random base station deployment, deriving tractable expressions for coverage probability and ergodic rate capturing the interplay between caching gains and NOMA's spectral efficiency. Guo et al. \cite{b_guo_edge} examined integration for 5G/6G networks, highlighting the need for intelligent cache placement aware of NOMA constraints and joint optimization of caching and power allocation.

Joint optimization has emerged as critical, recognizing independent optimization leaves synergistic gains unrealized. Cao et al. \cite{b_cao_joint} formulated joint caching and power allocation problems maximizing throughput while ensuring minimum rates for cell-edge users. Their iterative algorithm demonstrated significant improvements over separate optimization but acknowledged computational complexity from combinatorial cache placement coupled with non-convex power constraints. Liu et al. \cite{b_liu_uav} investigated UAV-assisted NOMA with caching for 6G, where flying base stations provide flexible coverage. Their analysis revealed UAV mobility enables dynamic cache positioning adapting to temporal traffic variations, while NOMA benefits from favorable line-of-sight propagation enhancing SIC performance.

Machine learning and deep reinforcement learning have been applied to address limitations of traditional reactive policies that cannot anticipate popularity shifts. Chen et al. \cite{b_chen_ml} provided an overview of ML approaches for wireless caching, covering supervised learning for popularity prediction and reinforcement learning frameworks learning optimal policies through environment interaction. Elghariani et al. \cite{b_elghariani_pred} investigated predictive caching using Long Short-Term Memory networks capturing temporal dependencies, achieving 15-20\% higher cache hit ratios compared to LRU and LFU policies under time-varying patterns.

Most relevant to this project, Li et al. \cite{b_li_drl} pioneered Deep Reinforcement Learning for dynamic caching in NOMA-based networks, formulating cache management as a Markov Decision Process where states encompass cache occupancy and channel quality, actions are caching decisions, and rewards combine cache hits with NOMA transmission success. Their Deep Q-Network agent with experience replay and target networks demonstrated 18-25\% performance improvements over static caching, particularly in high-density scenarios. However, their work assumed perfect SIC and ideal channel state information, leaving a gap under realistic impairments with SIC errors of 5-15\%, channel estimation noise, and hardware constraints.

Relay-assisted cooperative NOMA has been explored to improve cell-edge performance. Bepari et al. \cite{b_bepari_partial} analyzed partial cooperative strategies where near users forward decoded far-user information when direct transmission proves insufficient, reducing far-user outage probability by 30-40\% at modest near-user throughput cost. This suggests cache-aided NOMA may benefit from incorporating cooperative mechanisms triggered adaptively based on cache status and channel conditions.

Despite significant progress, existing literature exhibits notable limitations motivating the current research. Most analytical studies assume ideal physical layer conditions including perfect SIC and accurate CSI, neglecting practical impairments. Joint optimization frameworks often require extensive computational resources and may converge to suboptimal local equilibria. Machine learning approaches have been developed largely independently of physical layer considerations, with limited investigation of how caching policies should adapt to NOMA constraints and channel quality variations. Performance evaluation frequently relies on simplified channel models or static scenarios, with insufficient analysis under realistic wireless propagation including Rayleigh fading, path loss, user mobility, and time-varying interference.

The current project addresses these gaps by developing an integrated cache-aided NOMA framework employing a Double Dueling Deep Q-Network. Unlike prior works that rely on traditional $\mathcal{O}(N_{\text{files}})$ eviction action spaces, our framework introduces a novel binary action-space abstraction making DRL remarkably tractable for large, real-world content catalogs. Furthermore, our work explicitly differs by formulating the first CIC-incentivized DRL proactive caching policy with direct NOMA-layer state feedback mapping. The comprehensive simulation incorporates realistic NOMA physical layer modeling with configurable SIC imperfection factors, mixed fading (Rayleigh/Rician), distance-dependent path loss, and dynamic Zipf-distributed requests. Performance evaluation spans cache hit ratio, system sum-rate, CIC benefit rate, and outage probability, characterizing trade-offs between caching efficiency and physical layer quality of service. Comparative analysis against LRU, LFU, and Static Top-K baselines validates the DRL approach's advantages in dynamic environments, establishing a foundation for practical deployment of intelligent cache-aided NOMA systems in sixth-generation wireless networks.

\thispagestyle{plain}

%-------------------------------------------------------------------------------
%   CHAPTER-III DIVIDER PAGE
%   Matches IOT.pdf page 16 format
%   Page numbering continues from previous chapters
%-------------------------------------------------------------------------------

\clearpage

\begin{center}
    \vspace*{10cm}
    
    {\fontsize{16}{19.2}\selectfont\bfseries
    Chapter -- III}
    
    \vspace{1cm}
    
    {\fontsize{16}{19.2}\selectfont\bfseries
    PROBLEM DEFINITION AND\\
    MODULES}
    
    \vspace*{\fill}
\end{center}

\thispagestyle{plain}

%-------------------------------------------------------------------------------
%   CHAPTER 3: PROBLEM DEFINITION AND MODULES (With External Image)
%-------------------------------------------------------------------------------

\clearpage

\section{Problem Definition and Modules}

\subsection{Problem Definition}

Cache-Aided Non-Orthogonal Multiple Access (NOMA) emerges as a solution by combining edge caching with power-domain multiplexing. NOMA enables multiple users to share time-frequency resources through Successive Interference Cancellation (SIC), improving spectral efficiency by 30-50\%. Edge caching reduces latency from 50-200 ms (backhaul) to 1-5 ms (local cache) and offloads 40-60\% of traffic. However, integrating these technologies requires intelligent resource management in dynamic environments where user density, channel conditions, and content popularity fluctuate rapidly.

Traditional caching policies like Least Recently Used (LRU) and Least Frequently Used (LFU) operate reactively based on historical patterns, unable to anticipate future popularity shifts or adapt to physical layer constraints. They make decisions independently of NOMA transmission conditions, channel quality, and user distribution, missing optimization opportunities. Static policies like Top-K caching perform well under stable popularity but fail catastrophically when trends change.

\textbf{Core Research Problem:} Design and implement an intelligent cache-aided NOMA system using Deep Reinforcement Learning with a novel binary action space that learns optimal proactive caching policies. The system must adapt to dynamic network conditions, explicitly maximizing Cache-aided Interference Cancellation (CIC) opportunities, minimizing access latency, and ensuring acceptable quality of service under realistic wireless impairments including imperfect SIC and mixed fading channels.

\begin{figure}[htbp]
    \centering
    \includegraphics[width=0.55\textwidth]{system_model.png}
    \caption{Cache-Aided NOMA System Model showing base station with edge cache, near and far users with NOMA power allocation, and channel conditions}
    \label{fig:system_model}
\end{figure}

The proposed solution employs Deep Q-Networks (DQN) with dueling architecture, prioritized experience replay, and double Q-learning to learn optimal caching policies. The DQN agent observes system state including cache contents, recent request history, channel quality, and NOMA success rates, then selects actions to cache or evict content items to maximize long-term reward combining cache hits, transmission quality, and user fairness.

\subsubsection{System Model}
%====================================================

%====================================================
\subsubsection{Channel Model}
%====================================================

The wireless channel between the BS and each user is assumed to be frequency-flat within an individual Resource Block (RB). Let $h_{k,m}$ denote the complex channel coefficient between the BS and user $k$ over the $m$-th RB. The overall channel consists of three major propagation components, namely distance-dependent path loss, log-normal shadowing, and small-scale fading. Therefore, the instantaneous channel gain can be expressed as

\begin{equation}
|h_{k,m}|^2=
\frac{\xi_k |g_{k,m}|^2}
{PL_k},
\label{eq:channel_gain}
\end{equation}

where

\begin{itemize}
\item $PL_k$ denotes the large-scale path loss,
\item $\xi_k$ represents the shadowing coefficient,
\item $g_{k,m}$ denotes the small-scale fading coefficient.
\end{itemize}

The distance-dependent path loss is modeled as

\begin{equation}
PL_k=d_k^{\beta},
\label{eq:pathloss}
\end{equation}

where $d_k$ is the distance between the BS and user $k$, while $\beta$ is the path-loss exponent.

The shadowing component follows a log-normal distribution

\begin{equation}
10\log_{10}(\xi_k)\sim\mathcal{N}(0,\sigma_{sh}^2),
\end{equation}

where $\sigma_{sh}$ denotes the shadowing standard deviation.

The small-scale fading coefficient is assumed to follow Rayleigh fading, i.e.,

\begin{equation}
g_{k,m}\sim\mathcal{CN}(0,1),
\end{equation}

which results in exponentially distributed channel power gains.

The received signal at user $k$ over the $m$-th RB is

\begin{equation}
y_{k,m}
=
h_{k,m}x_m+n_{k,m},
\label{eq:received_signal}
\end{equation}

where $x_m$ denotes the transmitted superposition signal and

\begin{equation}
n_{k,m}\sim\mathcal{CN}(0,\sigma^2)
\end{equation}

is the additive white Gaussian noise (AWGN).

To facilitate NOMA transmission, users are arranged in ascending order according to their instantaneous channel gains

\begin{equation}
|h_{1,m}|^2
\le
|h_{2,m}|^2
\le
\cdots
\le
|h_{K,m}|^2.
\label{eq:ordering}
\end{equation}

For each RB, the user having the minimum channel gain is referred to as the weak user ($w$), whereas the user with the highest channel gain is referred to as the strong user ($s$). The proposed system adopts an extreme pairing strategy in which the strongest user is paired with the weakest user to maximize the channel gain difference and improve the effectiveness of power-domain multiplexing.

%====================================================
\subsubsection{Cache Placement Model}
%====================================================

The proposed framework adopts a proactive cache placement strategy in which frequently requested contents are stored at user devices during off-peak network hours. The cache placement phase is performed before the actual transmission process and therefore does not consume radio resources during data delivery.

Let the entire content library be represented by

\begin{equation}
\mathcal{F}
=
\{f_1,f_2,\cdots,f_N\},
\end{equation}

where each file has identical size.

Each user is equipped with a local cache memory of capacity

\begin{equation}
C<N,
\end{equation}

allowing only a subset of the library to be stored locally.

The popularity of the files follows the Zipf distribution

\begin{equation}
p_n=
\frac{n^{-\alpha}}
{\sum_{i=1}^{N}i^{-\alpha}},
\qquad
n=1,\cdots,N,
\label{eq:zipf_distribution}
\end{equation}

where $\alpha$ controls the popularity skewness.

During the cache placement phase, the BS stores the most popular contents at users according to the cache capacity constraint

\begin{equation}
\sum_{n=1}^{N}c_{k,n}\le C,
\label{eq:cache_constraint}
\end{equation}

where

\begin{equation}
c_{k,n}
=
\begin{cases}
1, & \text{if file }f_n\text{ is cached at user }k,\\
0, & \text{otherwise}.
\end{cases}
\label{eq:cache_indicator}
\end{equation}

When a user requests a file already available in its cache, the requested content is served locally without occupying wireless resources. Otherwise, the BS delivers the requested file through the downlink NOMA transmission.

Unlike conventional cache-enabled systems, the proposed Hybrid SIC/CIC framework also exploits cached information to suppress multi-user interference. Consider a NOMA pair consisting of a weak user ($w$) and a strong user ($s$). Let

\begin{equation}
c_w=
\begin{cases}
1,& \text{if the strong user's requested file is cached at the weak user},\\
0,& \text{otherwise},
\end{cases}
\end{equation}

and

\begin{equation}
c_s=
\begin{cases}
1,& \text{if the weak user's requested file is cached at the strong user},\\
0,& \text{otherwise}.
\end{cases}
\end{equation}

The binary variables $c_w$ and $c_s$ determine whether Cache-Aided Interference Cancellation (CIC) can be activated at the corresponding receiver. If $c_w=1$, the weak user can reconstruct and perfectly remove the strong user's interference directly from its local cache. Similarly, when $c_s=1$, the strong user can eliminate the weak user's signal without relying on successive interference cancellation.

Accordingly, four different transmission scenarios are possible:

\begin{enumerate}
\item $c_w=0,\;c_s=0$: Conventional NOMA with SIC.
\item $c_w=1,\;c_s=0$: Weak-user Cache-Aided Interference Cancellation.
\item $c_w=0,\;c_s=1$: Strong-user Cache-Aided Interference Cancellation.
\item $c_w=1,\;c_s=1$: Full Hybrid CIC, where both users cancel interference using cached information.
\end{enumerate}

This cache-aware transmission strategy enables the proposed system to dynamically switch between conventional SIC and CIC according to the cache status of each NOMA pair, thereby improving decoding reliability, reducing residual interference, and lowering the overall outage probability.

%====================================================
\subsubsection{User Pairing Strategy}
%====================================================

To improve spectral efficiency and exploit the power-domain multiplexing capability of NOMA, two users are grouped together on each Resource Block (RB). The proposed framework adopts an extreme user pairing strategy, where the user experiencing the strongest channel condition is paired with the user having the weakest channel condition within the same RB. Such pairing maximizes the channel gain difference between paired users, thereby improving power allocation efficiency and decoding performance.

Let the ordered channel gains over the $m$-th RB be

\begin{equation}
|h_{1,m}|^2
\le
|h_{2,m}|^2
\le
\cdots
\le
|h_{K,m}|^2.
\end{equation}

The paired users are denoted as

\begin{equation}
(w,s),
\end{equation}

where

\begin{equation}
|h_{w,m}|^2<|h_{s,m}|^2.
\end{equation}

The weak user is allocated a larger transmit power, while the strong user receives a smaller power allocation. Let

\begin{equation}
\alpha_{w,m}+\alpha_{s,m}=1,
\end{equation}

where

\begin{equation}
\alpha_{w,m}>\alpha_{s,m}.
\end{equation}

This power allocation guarantees that the weak user's Quality-of-Service (QoS) requirement is satisfied while allowing the strong user to decode the weak user's signal using successive interference cancellation (SIC) whenever cache-assisted cancellation is unavailable.

Each Resource Block serves only one NOMA pair. The RB allocation indicator is defined as

\begin{equation}
\beta_{k,m}
=
\begin{cases}
1,&\text{if user }k\text{ is assigned to RB }m,\\
0,&\text{otherwise}.
\end{cases}
\end{equation}

The RB assignment satisfies

\begin{equation}
\sum_{k=1}^{K}\beta_{k,m}\le2,
\qquad
\forall m,
\end{equation}

ensuring that at most two users share the same RB.

%====================================================
\subsubsection{Cache-Aided NOMA Transmission}
%====================================================

After user pairing, the BS simultaneously transmits the requested information to both users through power-domain superposition coding.

The transmitted signal over the $m$-th RB is

\begin{equation}
x_m
=
\sqrt{\alpha_{w,m}P_m}x_w
+
\sqrt{\alpha_{s,m}P_m}x_s,
\label{eq:superposition}
\end{equation}

where

\begin{itemize}
\item $P_m$ denotes the transmit power allocated to the $m$-th RB,
\item $x_w$ and $x_s$ represent the unit-power information symbols intended for the weak and strong users,
\item $\alpha_{w,m}$ and $\alpha_{s,m}$ denote the corresponding power allocation coefficients.
\end{itemize}

The received signal at user $k$ is

\begin{equation}
y_{k,m}
=
h_{k,m}
\left(
\sqrt{\alpha_{w,m}P_m}x_w
+
\sqrt{\alpha_{s,m}P_m}x_s
\right)
+n_{k,m},
\end{equation}

where

\begin{equation}
n_{k,m}\sim\mathcal{CN}(0,\sigma^2).
\end{equation}

Unlike conventional NOMA systems, the proposed framework utilizes cached contents not only to reduce backhaul traffic but also to suppress multi-user interference before signal detection.

Depending on the cache status of the paired users, the receiver dynamically switches between conventional SIC and Cache-Aided Interference Cancellation (CIC). Consequently, the received signal can be processed under four different operating conditions, thereby improving the overall decoding reliability and reducing transmission outages.

%====================================================
\subsubsection{Hybrid SIC/CIC Receiver Model}
%====================================================

In conventional NOMA systems, the weak user directly decodes its own signal by treating the strong user's transmission as interference. The strong user first decodes the weak user's signal and subsequently removes it using successive interference cancellation (SIC). However, practical SIC suffers from imperfect signal reconstruction, which leaves residual interference and degrades system performance.

To overcome this limitation, the proposed framework introduces a Hybrid SIC/CIC receiver that dynamically selects the interference cancellation mechanism according to the cache status of each user.

Let

\begin{equation}
c_w=
\begin{cases}
1,&\text{if the strong user's requested file is cached at the weak user},\\
0,&\text{otherwise},
\end{cases}
\end{equation}

and

\begin{equation}
c_s=
\begin{cases}
1,&\text{if the weak user's requested file is cached at the strong user},\\
0,&\text{otherwise}.
\end{cases}
\end{equation}

The variables $c_w$ and $c_s$ determine whether cache-assisted interference cancellation can be activated.

Accordingly, the receiver operates under one of the following four transmission modes.

\paragraph*{Mode 1: Conventional SIC}

When

\[
c_w=0,\qquad c_s=0,
\]

no useful cached information is available. Therefore, both users follow the conventional NOMA decoding procedure. The weak user treats the strong-user signal as interference, whereas the strong user performs imperfect SIC.

\paragraph*{Mode 2: Weak-User CIC}

When

\[
c_w=1,\qquad c_s=0,
\]

the weak user reconstructs the interfering signal directly from its local cache and removes it before decoding its desired signal. The strong user continues to employ conventional SIC.

\paragraph*{Mode 3: Strong-User CIC}

When

\[
c_w=0,\qquad c_s=1,
\]

the strong user already possesses the weak user's requested file in its cache. Therefore, it reconstructs the interference perfectly without performing SIC. The weak user still treats the strong-user transmission as interference.

\paragraph*{Mode 4: Full Hybrid CIC}

When

\[
c_w=1,\qquad c_s=1,
\]

both users have access to each other's requested content through local caching. Consequently, each receiver reconstructs and removes the interfering signal independently, eliminating the need for SIC completely.

The proposed Hybrid SIC/CIC receiver therefore provides a generalized interference management framework that seamlessly switches between conventional SIC and cache-assisted interference cancellation according to the instantaneous cache availability. This adaptive receiver improves decoding reliability, reduces residual interference, and forms the basis for the unified SINR and outage analysis presented in the subsequent section.
%====================================================
\subsubsection{Unified SINR Analysis}
%====================================================

The proposed Hybrid SIC/CIC receiver operates according to the cache availability of the paired users. Instead of deriving separate SINR expressions for each operating mode, a unified mathematical model is developed to represent all possible transmission scenarios.

Let

\begin{equation}
c_w,c_s\in\{0,1\},
\end{equation}

be binary cache indicators, where

\begin{itemize}
\item $c_w=1$ indicates that the weak user stores the strong user's requested content,
\item $c_s=1$ indicates that the strong user stores the weak user's requested content.
\end{itemize}

If the required interfering signal is available in the local cache, the corresponding interference can be perfectly reconstructed and removed before signal detection.

Therefore, the received SINR of the weak user can be expressed in a unified form as

\begin{equation}
\gamma_{w,m}
=
\frac{\alpha_{w,m}P_m|h_{w,m}|^2}
{(1-c_w)\alpha_{s,m}P_m|h_{w,m}|^2+\sigma^2}.
\label{eq:UnifiedWeakSINR}
\end{equation}

Similarly, the SINR of the strong user is

\begin{equation}
\gamma_{s,m}
=
\frac{\alpha_{s,m}P_m|h_{s,m}|^2}
{(1-c_s)\zeta\alpha_{w,m}P_m|h_{s,m}|^2+\sigma^2},
\label{eq:UnifiedStrongSINR}
\end{equation}

where $\zeta$ denotes the residual interference factor caused by imperfect SIC.

The above equations automatically represent all four receiver operating modes.

\textbf{Case 1: Conventional SIC}

When

\[
c_w=0,\qquad c_s=0,
\]

(\ref{eq:UnifiedWeakSINR}) and (\ref{eq:UnifiedStrongSINR}) become

\begin{equation}
\gamma_w=
\frac{\alpha_wP|h_w|^2}
{\alpha_sP|h_w|^2+\sigma^2},
\end{equation}

\begin{equation}
\gamma_s=
\frac{\alpha_sP|h_s|^2}
{\zeta\alpha_wP|h_s|^2+\sigma^2},
\end{equation}

which correspond to the conventional imperfect SIC receiver.

\textbf{Case 2: Weak-User CIC}

For

\[
c_w=1,\qquad c_s=0,
\]

the interference observed by the weak user is completely removed,

\begin{equation}
\gamma_w^{CIC}
=
\frac{\alpha_wP|h_w|^2}
{\sigma^2},
\end{equation}

whereas the strong user still performs conventional SIC.

\textbf{Case 3: Strong-User CIC}

When

\[
c_w=0,\qquad c_s=1,
\]

the strong user performs perfect interference cancellation using its cached content,

\begin{equation}
\gamma_s^{CIC}
=
\frac{\alpha_sP|h_s|^2}
{\sigma^2},
\end{equation}

while the weak user follows the conventional NOMA decoding procedure.

\textbf{Case 4: Full Hybrid CIC}

If

\[
c_w=1,\qquad c_s=1,
\]

both users eliminate interference using cached information and their SINRs become

\begin{equation}
\gamma_w
=
\frac{\alpha_wP|h_w|^2}
{\sigma^2},
\end{equation}

\begin{equation}
\gamma_s
=
\frac{\alpha_sP|h_s|^2}
{\sigma^2}.
\end{equation}

Hence, the proposed unified formulation provides a generalized SINR model for both conventional SIC and cache-assisted interference cancellation.

%====================================================
\subsubsection{Achievable Rate Analysis}
%====================================================

Based on Shannon's channel capacity theorem, the instantaneous achievable data rate of user $k$ over the $m$-th Resource Block is

\begin{equation}
R_{k,m}
=
B
\log_2
\left(
1+\gamma_{k,m}
\right),
\end{equation}

where $B$ denotes the bandwidth allocated to one RB.

Accordingly, the achievable rates of the weak and strong users are

\begin{equation}
R_w
=
B
\log_2
(1+\gamma_w),
\end{equation}

and

\begin{equation}
R_s
=
B
\log_2
(1+\gamma_s),
\end{equation}

respectively.

The achievable sum rate of one NOMA pair is

\begin{equation}
R_{pair}
=
R_w+R_s.
\end{equation}

For a network containing $L$ independent NOMA pairs, the overall system throughput is

\begin{equation}
R_{sum}
=
\sum_{i=1}^{L}
R_{pair,i}.
\end{equation}

Since the proposed Hybrid SIC/CIC receiver suppresses multi-user interference whenever cached information is available, both $\gamma_w$ and $\gamma_s$ increase. Consequently, the achievable data rate also improves compared with the conventional SIC-based receiver.

%====================================================
\subsubsection{Weak User Outage Probability}
%====================================================

Reliable communication is achieved only when the received SINR exceeds a predefined threshold. Let the minimum Quality-of-Service (QoS) requirement be denoted by the target transmission rate $R_{th}$.

The corresponding SINR threshold is

\begin{equation}
\gamma_{th}
=
2^{R_{th}}-1.
\end{equation}

An outage event occurs whenever the received SINR falls below this threshold.

Therefore, the outage probability of the weak user is defined as

\begin{equation}
P_{out,w}
=
Pr
\left(
\gamma_w
<
\gamma_{th}
\right).
\end{equation}

Substituting the unified SINR expression from (\ref{eq:UnifiedWeakSINR}),

\begin{equation}
P_{out,w}
=
Pr
\left(
\frac{\alpha_wP|h_w|^2}
{(1-c_w)\alpha_sP|h_w|^2+\sigma^2}
<
\gamma_{th}
\right).
\end{equation}

For the conventional SIC receiver,

\[
c_w=0,
\]

the outage probability becomes

\begin{equation}
P_{out,w}^{SIC}
=
Pr
\left(
\frac{\alpha_wP|h_w|^2}
{\alpha_sP|h_w|^2+\sigma^2}
<
\gamma_{th}
\right).
\end{equation}

When cache-aided interference cancellation is activated,

\[
c_w=1,
\]

the interference term disappears and the outage probability reduces to

\begin{equation}
P_{out,w}^{CIC}
=
Pr
\left(
\frac{\alpha_wP|h_w|^2}
{\sigma^2}
<
\gamma_{th}
\right).
\end{equation}

Since the denominator contains only thermal noise, the weak user experiences a considerably lower outage probability under the proposed Hybrid SIC/CIC receiver than under the conventional SIC-based NOMA system.
%====================================================
\subsubsection{Strong User Outage Probability}
%====================================================

The strong user first removes the weak user's signal before decoding its own information. In the conventional NOMA receiver, this process is performed using Successive Interference Cancellation (SIC). However, practical SIC is not perfect and leaves a small amount of residual interference due to channel estimation errors and imperfect signal reconstruction.

Considering imperfect SIC, the outage probability of the strong user is defined as

\begin{equation}
P_{out,s}
=
\Pr
\left(
\gamma_s
<
\gamma_{th}
\right),
\end{equation}

where $\gamma_{th}=2^{R_{th}}-1$ denotes the target SINR threshold.

Substituting the unified SINR expression in (\ref{eq:UnifiedStrongSINR}), the outage probability becomes

\begin{equation}
P_{out,s}
=
\Pr
\left(
\frac{\alpha_sP|h_s|^2}
{(1-c_s)\zeta\alpha_wP|h_s|^2+\sigma^2}
<
\gamma_{th}
\right).
\end{equation}

Depending on the cache status of the strong user, two different decoding cases are obtained.

\paragraph*{Case 1: Conventional SIC}

When the strong user does not store the weak user's requested content, i.e.,

\[
c_s=0,
\]

the receiver employs imperfect SIC and the outage probability is

\begin{equation}
P_{out,s}^{SIC}
=
\Pr
\left(
\frac{\alpha_sP|h_s|^2}
{\zeta\alpha_wP|h_s|^2+\sigma^2}
<
\gamma_{th}
\right).
\end{equation}

The residual interference coefficient $\zeta$ limits the achievable SINR and therefore increases the outage probability.

\paragraph*{Case 2: Cache-Aided Interference Cancellation}

When the strong user already stores the weak user's requested content,

\[
c_s=1,
\]

the interfering signal is reconstructed directly from the cache without performing SIC. Consequently, the residual interference is completely eliminated and the received SINR becomes

\begin{equation}
\gamma_s^{CIC}
=
\frac{\alpha_sP|h_s|^2}
{\sigma^2}.
\end{equation}

The corresponding outage probability is

\begin{equation}
P_{out,s}^{CIC}
=
\Pr
\left(
\frac{\alpha_sP|h_s|^2}
{\sigma^2}
<
\gamma_{th}
\right).
\end{equation}

Since the residual interference term is completely removed, the proposed Hybrid SIC/CIC receiver significantly improves the decoding reliability of the strong user compared with the conventional SIC-based receiver.

%====================================================
\subsubsection{Pair Outage Probability}
%====================================================

A successful NOMA transmission requires that both the weak and the strong users correctly decode their intended signals. Therefore, a transmission outage occurs whenever at least one user experiences an outage.

Let

\begin{equation}
P_{out,w}
\end{equation}

and

\begin{equation}
P_{out,s}
\end{equation}

denote the outage probabilities of the weak and strong users, respectively.

Assuming independent decoding events for the paired users, the successful transmission probability of one NOMA pair is

\begin{equation}
P_{success}
=
(1-P_{out,w})
(1-P_{out,s}).
\end{equation}

Accordingly, the outage probability of the NOMA pair is

\begin{equation}
P_{pair}
=
1-
(1-P_{out,w})
(1-P_{out,s}).
\label{eq:pairout}
\end{equation}

Equation (\ref{eq:pairout}) indicates that the outage probability decreases whenever either user's decoding reliability improves.

For the proposed Hybrid SIC/CIC framework, the pair outage probability automatically adapts to the cache status of both users through the unified SINR expressions. Consequently,

\begin{itemize}

\item when neither user possesses cached information ($c_w=0$, $c_s=0$), (\ref{eq:pairout}) reduces to the conventional SIC-NOMA pair outage;

\item when only the weak user performs cache-aided interference cancellation ($c_w=1$, $c_s=0$), the weak-user outage decreases while the strong-user outage remains unchanged;

\item when only the strong user performs cache-aided interference cancellation ($c_w=0$, $c_s=1$), the strong-user outage decreases;

\item when both users exploit locally cached contents ($c_w=1$, $c_s=1$), the outage probabilities of both users are minimized, resulting in the lowest pair outage probability.

\end{itemize}

Therefore, the proposed Hybrid SIC/CIC receiver improves the reliability of every NOMA pair by dynamically selecting the appropriate interference cancellation mechanism according to the instantaneous cache availability.
%====================================================
\subsubsection{End-to-End System Outage Probability}
%====================================================

While the pair outage probability evaluates the reliability of an individual NOMA pair, the overall network performance depends on the successful transmission of all scheduled NOMA pairs. Therefore, the end-to-end outage probability is adopted as the primary reliability metric of the proposed Hybrid SIC/CIC framework.

Assume that the network contains $L$ independent NOMA pairs. Let
$P_{pair,i}$ denote the outage probability of the $i$-th pair.

The successful transmission probability of the entire system is

\begin{equation}
P_{success}^{sys}
=
\prod_{i=1}^{L}
\left(1-P_{pair,i}\right).
\label{eq:success_system}
\end{equation}

Consequently, the end-to-end outage probability is obtained as

\begin{equation}
P_{E2E}
=
1-
\prod_{i=1}^{L}
\left(1-P_{pair,i}\right).
\label{eq:e2e}
\end{equation}

For a sufficiently large number of user pairs, the average system outage can alternatively be expressed as

\begin{equation}
P_{sys}
=
\frac{1}{L}
\sum_{i=1}^{L}
P_{pair,i}.
\label{eq:avg_system}
\end{equation}

Equation (\ref{eq:e2e}) demonstrates that the overall network reliability depends on the outage performance of every scheduled NOMA pair. Since the proposed Hybrid SIC/CIC receiver reduces the outage probability of both weak and strong users through cache-assisted interference cancellation, the end-to-end outage probability is significantly lower than that of the conventional SIC-based NOMA system.

%====================================================
\subsubsection{Dynamic Cache-Aware Power Allocation}
%====================================================

Unlike conventional NOMA systems that employ fixed power allocation coefficients, the proposed framework dynamically adjusts the transmit power according to the instantaneous channel conditions and cache availability of each NOMA pair.

Let

\begin{equation}
\alpha_{w,m}+\alpha_{s,m}=1,
\label{eq:power_sum}
\end{equation}

where

\begin{equation}
0<\alpha_{s,m}<\alpha_{w,m}<1.
\end{equation}

The power allocation coefficients are determined according to the receiver operating mode.

\textbf{Mode 1: Conventional SIC}

When no cache-assisted interference cancellation is available,

\[
c_w=0,\qquad c_s=0,
\]

the weak user requires higher transmit power to overcome multi-user interference. Therefore,

\begin{equation}
\alpha_{w,m}\gg\alpha_{s,m}.
\end{equation}

This corresponds to the conventional NOMA power allocation strategy.

\textbf{Mode 2: Weak-User CIC}

If

\[
c_w=1,
\]

the weak user removes the interference using its cached information. Hence, the minimum power required to satisfy the target SINR becomes

\begin{equation}
\alpha_{w,m}^{min}
=
\frac{\gamma_{th}\sigma^2}
{P_m|h_{w,m}|^2}.
\label{eq:weakpower}
\end{equation}

The remaining transmit power can therefore be allocated to the strong user, improving the overall system throughput.

\textbf{Mode 3: Strong-User CIC}

When

\[
c_s=1,
\]

the residual SIC interference is completely eliminated. Consequently, the feasible power region expands as

\begin{equation}
\alpha_{w,m}^{max}
=
1-
\frac{\gamma_{th}\sigma^2}
{P_m|h_{s,m}|^2}.
\label{eq:strongpower}
\end{equation}

This additional degree of freedom allows more efficient power sharing between the paired users.

\textbf{Mode 4: Full Hybrid CIC}

When both users possess the interfering content,

\[
c_w=c_s=1,
\]

the proposed algorithm jointly determines the optimal power coefficients

\[
(\alpha_w,\alpha_s),
\]

such that the outage probability is minimized while maintaining the required Quality-of-Service constraints for both users.

Consequently, the proposed cache-aware power allocation dynamically adapts the transmit power according to both channel quality and cache status, leading to improved spectral efficiency, lower outage probability, and higher energy efficiency.

%====================================================
\subsubsection{Optimization Problem}
%====================================================

The objective of the proposed Hybrid SIC/CIC framework is to jointly optimize the power allocation coefficients for all NOMA pairs while exploiting cache-assisted interference cancellation to improve network reliability.

The optimization problem is formulated as

\begin{equation}
\begin{aligned}
\max_{\{\alpha_{w,m},\alpha_{s,m}\}}
\quad
&
\sum_{m=1}^{M}
\left(
R_{w,m}+R_{s,m}
\right)
\\
\text{s.t.}
\quad
&
P_{out,w}\le\epsilon_w,
\\
&
P_{out,s}\le\epsilon_s,
\\
&
R_{w,m}\ge R_{th},
\\
&
R_{s,m}\ge R_{th},
\\
&
\alpha_{w,m}+\alpha_{s,m}=1,
\\
&
0<\alpha_{s,m}<\alpha_{w,m}<1,
\\
&
P_m\le P_{max},
\\
&
c_w,c_s\in\{0,1\},
\end{aligned}
\label{eq:optimization}
\end{equation}

where

\begin{itemize}

\item $R_{th}$ denotes the minimum QoS data rate,

\item $\epsilon_w$ and $\epsilon_s$ represent the maximum allowable outage probabilities,

\item $P_{max}$ is the BS transmit power budget,

\item $c_w$ and $c_s$ indicate the cache status of the paired users.

\end{itemize}

The optimization problem in (\ref{eq:optimization}) is non-convex because the achievable rates and outage probabilities are nonlinear functions of the power allocation coefficients and the cache-assisted interference cancellation states. Therefore, instead of relying on computationally expensive iterative optimization methods, this work employs a Deep Q-Network (DQN)-based learning framework to determine near-optimal power allocation decisions. The DQN agent continuously observes the network state, including channel gains, cache availability, and user QoS requirements, and learns an adaptive power allocation policy that maximizes the long-term system sum rate while simultaneously minimizing the end-to-end outage probability under the given transmit power constraints.


%====================================================
\subsubsection{Achievable Rate}
%====================================================

Let

\begin{equation}
\rho=\frac{P_m}{B N_0}
\end{equation}

denote the transmit SNR.

The achievable rate of the weak user is

\begin{equation}
R_{w,m}
=
B\log_2
\left(
1+\gamma_{w,m}
\right),
\end{equation}

while the strong user's rate is

\begin{equation}
R_{s,m}
=
B\log_2
\left(
1+\gamma_{s,m}
\right).
\end{equation}

The total NOMA throughput is

\begin{equation}
R_{\mathrm{sum}}
=
\sum_{m=1}^{M}
\left(
R_{w,m}+R_{s,m}
\right).
\end{equation}

%====================================================
\subsubsection{Resource Block Allocation}
%====================================================

Let

\begin{equation}
\beta_{k,m}
=
\begin{cases}
1,& \text{if RB }m\text{ is assigned to user }k,\\
0,& \text{otherwise}.
\end{cases}
\end{equation}

Each RB can serve at most two NOMA users:

\begin{equation}
\sum_{k=1}^{K}
\beta_{k,m}
\le 2,
\qquad \forall m.
\end{equation}

The total power allocation satisfies

\begin{equation}
\sum_{k=1}^{K}
\sum_{m=1}^{M}
\alpha_{k,m}
\le 1.
\end{equation}


% IV. PROPOSED DRL CACHING FRAMEWORK
%=============================================================
\subsubsection{Proposed DRL Caching Framework}

To effectively outline the online placement boundaries applied immediately following the Hybrid Cache-Aided TDMA-NOMA transmission layer outlined in the System Model, the formalized execution parameters are mapped strictly within Algorithm \ref{alg:dqn}.

\begin{algorithm}
\caption{Hybrid Cache-Aware TDMA-NOMA Double Dueling DQN Cache Training}
\label{alg:dqn}
\begin{algorithmic}[1]
\REQUIRE User set $K=200$, File catalog $N=2000$, Cache capacity $C=200$, Maximum training episodes $E = 500$
\STATE \textbf{Variables:} Online network $\theta$ (the primary model evaluating real-time actions), Target network $\theta^-$ (delayed clone for stability), Replay buffer $\mathcal{B}$
\STATE Initialize online network $\theta$, target network $\theta^- \leftarrow \theta$, and bounded replay buffer $\mathcal{B}$
\STATE Assign exploration factor $\epsilon \leftarrow 1.0$, global timeline $t \leftarrow 0$
\FOR{episode $e = 1$ to $E = 500$}
    \STATE Reset cache buffer; instantiate fresh user positions, channel conditions, and cache states
    \STATE Generate random request stream $(u_i, f_i)_{i=1}^{T}$ for $T=10000$
    \FOR{each dynamic request $(u_i, f_i)$}
        \STATE Extract feature state $s_t \in \mathbb{R}^{14}$
        \STATE With probability $\epsilon$: $a_t \leftarrow$ Random Action; else $a_t \leftarrow \arg\max_a Q(s_t, a; \theta)$
        \IF{$a_t = 1$ and target file absent from cache}
            \STATE Insert file $f_i$; eject least popular entry if memory full
        \ENDIF
        \STATE Sort active users according to channel gain
        \STATE Form strong--weak user pairs using channel disparity and cache availability
        \STATE Allocate each selected pair to the corresponding TDMA transmission slot
        \STATE Perform Hybrid TDMA-NOMA superposition transmission
        \STATE Apply Cache-Aided Interference Cancellation (CIC) followed by SIC whenever required
        \STATE Compile cache-hit ratio, CIC success, achievable sum-rate and energy-efficiency metrics; observe reward $r_t$ via \eqref{eq:reward}
        \STATE Queue transition $(s_t, a_t, r_t, s_{t+1})$ into $\mathcal{B}$ with max priority flag
        \IF{$t > t_{\text{warmup}}$ and $(t \bmod 4 = 0)$}
            \STATE Extract IS-weighted transition minibatch from $\mathcal{B}$
            \STATE Calculate Double-Q target $y_t$ leveraging \eqref{eq:ddqn}
            \STATE Update bounds $\theta$ via Huber Loss backpropagation
            \STATE Soft-blend target network $\theta^-$ using $\tau = 0.005$
        \ENDIF
        \STATE Decay $\epsilon$ linearly; advance global tick $t \leftarrow t + 1$
    \ENDFOR
\ENDFOR
\end{algorithmic}
\end{algorithm}

\subsubsection{MDP Formulation}
We formulate the online cache placement problem as a continuous Markov Decision Process (MDP) defined by the tuple $\langle \mathcal{S}, \mathcal{A}, \mathcal{R}, \mathcal{P}, \gamma \rangle$.

\textbf{State space:} Upon receiving a request at step $t$, the extracted state vector $s_t \in \mathbb{R}^{14}$ encodes:
(i) Normalized popularity rank of the requested file,
(ii) Localized recent request frequency scalar of the target file,
(iii) Binary cache membership boolean indicator,
(iv) Cache storage occupancy ratio,
(v--vii) Scalar mean, min, and max historical popularity traits of currently cached files,
(viii) Binary quality-improvement boolean indicator for the present cache swap,
(ix) Popularity ratio relative to the absolute least popular cached file,
(x) Moving-average CIC potential rate from recent NOMA pairing history,
(xi--xii) Mean and standard deviation distributions of real-time channel gains,
(xiii) Historical NOMA transmission success tracking rate, and
(xiv) Evaluated file request diversity ratio bounded within the moving window.

\textbf{Action space:} In contrast to traditional policies evaluating $\mathcal{O}(N)$ eviction slots, our action definition is strictly binary: $a_t \in \{0, 1\}$. $a_t = 1$ forces the cache to encode the newly requested file (triggering an internal heuristic lowest-popularity eviction if full), whereas $a_t = 0$ actively skips the caching stage. This reduces learning complexity to $\mathcal{O}(1)$.

\textbf{Reward function:} The immediate state-action transition reward $r_t$ isolates strict decision boundary qualities:

\begin{equation}
r_t =
\begin{cases}
+2.0 & \text{cache hit (correct prior decision)} \\
+1.5 & \text{cache top-10\% popular file} \\
+0.8 & \text{cache top-30\% popular file} \\
+0.2 & \text{cache top-50\% popular file} \\
-0.5 & \text{cache bottom-50\% (unpopular file)} \\
+1.5 & \textbf{CIC enabled by this placement} \\
-0.5 & \text{skip a popular file with free space} \\
+0.1 & \text{correctly skip unpopular file}
\end{cases}
\label{eq:reward}
\end{equation}

Furthermore, explicit eviction quality bonuses ($\pm 0.5$) are applied exclusively when the storage bounds actuate a deletion (i.e., verifying if the swapped file historically underperformed the insertion candidate).

\subsubsection{Double Dueling DQN Architecture}
The central Q-function estimator relies on a \textit{Dueling DQN} topology \cite{b_wang_dueling} which intelligently segregates the raw state value $V(s)$ from action-specific advantages $A(s, a)$:

\begin{equation}
Q(s, a; \theta) = V(s; \theta_V) + \left[ A(s, a; \theta_A) - \frac{1}{|\mathcal{A}|}\sum_{a'} A(s, a'; \theta_A) \right]
\label{eq:dueling}
\end{equation}

The structural network dictates: Input($14$) $\rightarrow$ FC($64$)-ReLU $\rightarrow$ FC($32$)-ReLU. This trunk subsequently bifurcates into an autonomous value stream FC($16$)-ReLU-FC($1$) and an equivalent advantage stream FC($16$)-ReLU-FC($2$). To mitigate Q-value overestimation phenomena, the \textit{Double DQN} paradigm \cite{b_van_hasselt} decouples target action selection from mathematical evaluation:

\begin{equation}
y_t = r_t + \gamma \cdot Q_{\text{target}}\!\left(s_{t+1}, \arg\max_{a} Q_{\text{online}}(s_{t+1}, a)\right)
\label{eq:ddqn}
\end{equation}
enforcing an infinite horizon discount factor $\gamma = 0.99$.

\subsubsection{Prioritized Experience Replay}
State transitions are pushed into an asynchronous prioritized replay buffer \cite{b_schaul_per} comprising $50,000$ slots with a sampling exponent $\alpha_{\text{PER}} = 0.6$. Importance sampling (IS) weights equipped with $\beta$ annealed steadily from $0.4$ to $1.0$ over $5\times 10^6$ training steps adjust for non-uniform loss updates.

\subsubsection{Computational Complexity}
A critical advantage of the proposed framework is its runtime efficiency. While optimal proactive caching requires an NP-hard combinatorial search over $\mathcal{O}\left(\binom{N}{C}\right)$ subsets, the trained DRL agent operates in constant continuous time. The forward pass of the highly compact $14 \rightarrow 64 \rightarrow 32 \rightarrow 16 \rightarrow 2$ network involves merely $\approx 3{,}700$ FLOPs per request, achieving $\mathcal{O}(1)$ execution complexity. This ensures strict adherence to the ultra-low latency constraints demanded by 6G URLLC (Ultra-Reliable Low-Latency Communications) standards.


%=============================================================

\subsection{Modules}

The Cache-Aided NOMA system is decomposed into five integrated modules implementing physical layer simulation, caching management, deep reinforcement learning, content request generation, and performance evaluation.

\subsubsection{Module 1: NOMA Physical Layer Simulator}

The NOMA Physical Layer Simulator implements realistic downlink transmission modeling including channel propagation, power allocation, and Successive Interference Cancellation with practical impairments.

\textbf{Channel Modeling:} The simulator accurately tracks the mathematical definitions of path loss (exponent $\beta=3.5$) and Rayleigh fading to generate composite channel gains.

\textbf{NOMA Power Allocation:} The system implements a hybrid cache-aware dynamic power allocation strategy. The initial power coefficients are set to $(\alpha_w, \alpha_s) = (0.8, 0.2)$, but when cache status is available, the allocator adaptively adjusts coefficients based on CIC/SIC feasibility: if the weak user has cached the strong user's content ($c_w = 1$), the weak user experiences no interference and requires less power; if the strong user has cached content ($c_s = 1$), perfect SIC ($\zeta = 0$) is achieved, allowing more power to be redirected to the weak user. This is implemented via closed-form analytical bounds derived from the SINR constraints.

\textbf{SIC Implementation:} The receiver physically simulates the Hybrid SIC/CIC logic, transitioning seamlessly between conventional imperfect SIC (with $\zeta = 0.05$) and interference-free cache-aided reception based directly on the proactive caching state.

\textbf{Performance Metrics:} Calculates real-time system sum-rate and validates outage probability thresholds by checking $R_i < R_{\text{target}}$ corresponding to an SNR threshold of $0.3$ bps/Hz.

\subsubsection{Module 2: Cache Management System}

The Cache Management System implements multiple caching policies including reactive baselines and the proposed DQN-based adaptive policy.

\textbf{Cache Data Structure:} The cache is implemented as a hash map with ordered dictionary for efficient $O(1)$ lookup, insertion, and deletion. Cache capacity is $C = 200$ files from a catalog of $N = 2000$ files. Each cache entry stores file ID, size, timestamp, access count, and popularity score.

\textbf{Baseline Policies:} Least Recently Used (LRU) maintains access timestamps and evicts the content with oldest last-access time. Least Frequently Used (LFU) tracks cumulative request counts and removes the item with minimum count, with ties broken by recency. Static Top-K pre-caches the 200 most popular files based on Zipf distribution and never updates. First-In-First-Out (FIFO) evicts the oldest cached item regardless of usage patterns.

\textbf{Cache Operations:} The lookup operation checks if requested file ID exists in cache hash map. On cache hit, the file is retrieved locally with latency $T_{\text{cache}} = 2$ ms, access timestamp is updated, and counters are incremented. On cache miss, the file is fetched from backhaul with latency $T_{\text{backhaul}} = 100$ ms, and cache replacement policy determines which item to evict if cache is full.

\textbf{Performance Tracking:} The module maintains running statistics including total requests, cache hits, cache misses, cumulative access time, and per-file access counters. Cache hit ratio is computed as $\text{Hit Ratio} = \frac{\text{Hits}}{\text{Total Requests}}$. Average access time is $\text{AAT} = \frac{1}{n}\sum_{i=1}^n T_i$ where $T_i$ is access time for request $i$.

\subsubsection{Module 3: Deep Q-Network Agent}

The Deep Q-Network Agent implements the core machine learning component for intelligent caching decisions.

\textbf{State Representation:} The state vector $s_t$ is a novel NOMA-aware 14-dimensional representation that jointly encodes: (1) Cache occupancy features (hit rate, utilization); (2) Content popularity (request rate, unique requests); (3) Channel statistics (average channel gain, variance); and (4) NOMA pairing history (CIC opportunities, successful SIC rate). This highly condensed state avoids the curse of dimensionality associated with tracking individual file states.

\textbf{Action Space:} The action space employs a novel binary abstraction $\mathcal{A} \in \{0, 1\}$, where $a=1$ indicates caching the currently requested file (and evicting the least valuable item based on a heuristic if full), and $a=0$ indicates bypassing the cache. This reduces the action space from $\mathcal{O}(N_{\text{files}})$ to $\mathcal{O}(1)$, making DRL remarkably tractable for large, real-world content catalogs.

\textbf{DQN Architecture:} The Double Dueling DQN architecture splits the Q-network into a Value stream and an Advantage stream. The input state passes through fully connected layers (64 and 32 neurons, ReLU) before splitting into: a value stream FC(16)-ReLU-FC(1) and an advantage stream FC(16)-ReLU-FC(2). This architecture enables the network to learn which states are valuable independently of action choices.

\textbf{Reward Function:} The reward signal is explicitly CIC-incentivized. Instead of a continuous function, it relies on a discrete boundary evaluation as formulated in \eqref{eq:reward}, assigning high positive reinforcements (+2.0) for cache hits and (+1.5) for enabling multi-user NOMA cancellation.

\subsubsection{Module 4: Content Request Generator}

The Content Request Generator produces realistic request patterns with temporal dynamics.

\textbf{Zipf Distribution:} Base content popularity follows the Zipf law with probability $P(k) = \frac{k^{-\alpha}}{H_{N,\alpha}}$ where $H_{N,\alpha} = \sum_{i=1}^N i^{-\alpha}$ is the normalization constant. The skewness parameter $\alpha = 1.0$ creates a strong popularity concentration. The generated requests are explicitly bounded by the catalog limit $N=2000$.

\textbf{User Request Pattern:} Each of the $K=200$ active users generates requests independently sampled from the overarching Zipf distribution. The temporal request stream is aggregated at the base station to form the training and evaluation payload for the DQN agent.

\textbf{Request Metadata:} Each generated request encapsulates the file ID, the target user, and the instantaneous channel quality (Rayleigh fading gain). This metadata enables the DQN agent to make context-aware caching decisions considering both content popularity and transmission conditions.

\subsubsection{Module 5: Performance Evaluation and Visualization}

The Performance Evaluation module collects, aggregates, and visualizes comprehensive system metrics.

\textbf{Metric Collection:} Real-time collectors track cache hit ratio, average access time, NOMA success rate, outage probability, system sum-rate, and per-user rates. Metrics are computed over sliding windows of 100 requests for smoothing. Statistical aggregators maintain running mean, variance, minimum, maximum, and percentiles (5th, 50th, 95th) for all metrics.

\textbf{Comparative Analysis:} The module runs multiple policies (Static Top-K, LRU, LFU, Random, and the trained Double Dueling DQN) under identical conditions using the same random seed for request generation and channel realizations over 1000 Monte Carlo realizations. The framework evaluates performance across a dynamic SNR sweep from -10 dB to 30 dB.

\textbf{Visualization:} Performance plots include cache hit ratio, outage probability versus SNR curves, sum-rate evaluations, weak and strong user specific data rates, Bit Error Rate (BER), and energy efficiency measurements.

\textbf{Result Aggregation:} Experimental results are formatted into comprehensive summary outputs, tracking performance benchmarks at peak SNR levels (30 dB) allowing direct tabular verification against optimal thresholds.

\thispagestyle{plain}

%-------------------------------------------------------------------------------
%   CHAPTER-IV DIVIDER PAGE
%   Matches IOT.pdf page 23 format
%   Page numbering continues from previous chapters
%-------------------------------------------------------------------------------

\clearpage

\begin{center}
    \vspace*{10cm}
    
    {\fontsize{16}{19.2}\selectfont\bfseries
    Chapter -- IV}
    
    \vspace{1cm}
    
    {\fontsize{16}{19.2}\selectfont\bfseries
    SOFTWARE DESIGN}
    
    \vspace*{\fill}
\end{center}

\thispagestyle{plain}

%-------------------------------------------------------------------------------
%   CHAPTER 4: SOFTWARE DESIGN
%   Concise 3-page version matching IOT.pdf structure and created diagrams
%-------------------------------------------------------------------------------

\clearpage

\section{Software Design}

This chapter presents the software architecture and design of the cache-aided NOMA simulation framework. The design follows modular principles ensuring maintainability, extensibility, and clear separation of concerns between physical layer operations, caching management, and deep reinforcement learning components.

\subsection{System Architecture Overview}

The simulation framework adopts a layered architecture with four distinct layers, each encapsulating specific functionalities. Figure \ref{fig:system_architecture} illustrates the high-level system architecture showing interactions between major components.

\begin{figure}[htbp]
    \centering
    \includegraphics[width=0.85\textwidth]{system_architecture.png}
    \caption{System Architecture showing layered design with Application Layer, ML/DRL Layer, Caching Layer, and Physical Layer}
    \label{fig:system_architecture}
\end{figure}

\textbf{Application Layer:} The top layer manages request generation using Zipf distribution with skewness parameter $\alpha = 1.0$, performance monitoring collecting cache hit ratio, average access time, SINR, and outage probability, and result visualization using matplotlib for comparative plots and training convergence graphs.

\textbf{ML/DRL Layer:} Implements the Deep Q-Network agent with Dueling DQN architecture containing Q-network ($\theta$) and target network ($\theta^-$). The state encoder processes cache utilization, popularity metrics, channel statistics, and NOMA pairing history into a compact 14-dimensional NOMA-aware state vector. The agent utilizes a novel binary action space $\mathcal{A} \in \{0, 1\}$. Experience replay buffer stores transitions with prioritized sampling. The reward function explicitly incentivizes CIC opportunities using weighted sum $r_t = w_{hit} \cdot \mathbb{1}_{\text{hit}} + w_{cic} \cdot \mathbb{1}_{\text{cic}} + w_{rate} \cdot R_{sum,t} - w_{outage} \cdot \mathbb{1}_{\text{outage}}$.

\textbf{Caching Layer:} Manages content storage with hash map enabling $O(1)$ lookup for 200-file capacity from 2000-file catalog. Implements multiple replacement policies including LRU (least recently used), LFU (least frequently used), Static Top-K (pre-caching popular content), and DQN-based intelligent caching. Cache hits retrieve in 2ms while misses fetch from backhaul in 100ms.

\textbf{Physical Layer:} Implements NOMA transmission with Rayleigh fading channel model, path loss $L(d) = 120 + 35\log_{10}(d/100)$ dB, and 8 dB shadow fading. The system leverages a hybrid cache-aware power allocation strategy that dynamically adapts coefficients from a baseline $(0.8, 0.2)$ based on available CIC and perfect SIC ($\zeta=0$) opportunities. The SIC module models imperfect cancellation with factor $\zeta = 0.05$ (5\% residual interference) for non-cached payloads. SINR calculations for near user after SIC: $\gamma_2 = \frac{P_2|h_2|^2}{\zeta P_1|h_2|^2 + \sigma^2}$ and far user direct decoding: $\gamma_1 = \frac{P_1|h_1|^2}{P_2|h_1|^2 + \sigma^2}$.

\subsection{System Flow Diagram}

Figure \ref{fig:system_flow} presents the complete end-to-end workflow from initialization through request processing to performance analysis.

\begin{figure}[htbp]
    \centering
    \includegraphics[width=0.65\textwidth]{system_flow.png}
    \caption{Complete Cache-Aided NOMA System Flow diagram}
    \label{fig:system_flow}
\end{figure}

\textbf{Initialization Phase:} Loads configuration parameters including cache size (200), library size (2000), Zipf $\alpha = 1.0$, NOMA parameters ($\zeta = 0.05$, baseline power allocation 0.8/0.2), and DQN hyperparameters (learning rate $10^{-3}$, discount $\gamma = 0.99$, buffer size 50000, target update $\tau = 0.005$). Initializes Q-network with Double Dueling architecture (64-32 neurons) and target network with copied weights. Positions users at 400m (far) and 100m (near) from base station.

\textbf{Warm-up Phase:} Generates 500 initial requests to populate cache using baseline policy and collect training data for DQN, preventing cold-start problems.

\textbf{Main Simulation Loop:} Iterates for 50 episodes with 1000 requests each. For each request: (1) Generates file ID from Zipf distribution, (2) Observes state (cache status, request history, SINR), (3) Performs cache lookup - if hit retrieves in 2ms, if miss fetches from backhaul in 100ms and triggers replacement decision using DQN $\epsilon$-greedy action selection or baseline policy, (4) Executes NOMA transmission with power allocation and SIC decoding, (5) Calculates reward based on hit, latency, SINR, and outage, (6) Stores transition in replay buffer, (7) Samples mini-batch (64 transitions) with prioritized sampling if buffer size $\geq$ 64, (8) Updates Q-network minimizing TD error $\delta = (r + 0.99 \max_{a'} Q(s',a';\theta^-)) - Q(s,a;\theta)$, (9) Soft updates target network $\theta^- \leftarrow 0.005\theta + 0.995\theta^-$, (10) Decays exploration $\epsilon \leftarrow \max(0.01, 0.995\epsilon)$.

\textbf{Analysis Phase:} Aggregates statistics computing mean and 95\% confidence intervals across 50 episodes, performs comparative analysis between LRU, LFU, Static Top-K and DQN policies, generates performance plots showing cache hit ratio, outage probability, access time, and sum-rate, and exports results to CSV files.

\subsection{Detailed Process Flows}

\subsubsection{NOMA Transmission Process}

Figure \ref{fig:noma_flow} details the NOMA physical layer transmission procedure.

\begin{figure}[htbp]
    \centering
    \includegraphics[width=0.85\textwidth]{noma_transmission.png}
    \caption{NOMA Physical Layer Transmission Process}
    \label{fig:noma_flow}
\end{figure}

When request arrives, the system selects user pair (near at 100m, far at 400m), generates channel realizations using Rayleigh fading $h_i \sim \mathcal{CN}(0, \sigma_h^2)$ with path loss and shadowing, computes channel gains satisfying $|h_1|^2 < |h_2|^2$, adaptively allocates power dynamically evaluating CIC possibilities (or defaulting to $P_1 = 0.8P_{\text{total}}$ and $P_2 = 0.2P_{\text{total}}$), superposes signals at base station as $x(t) = \sqrt{P_1}s_1(t) + \sqrt{P_2}s_2(t)$, and transmits to both users. 

Near user performs three-step SIC: (1) Decodes far user signal treating own signal as interference with $\gamma_{2\rightarrow1} = \frac{P_1|h_2|^2}{P_2|h_2|^2 + \sigma^2}$, (2) Cancels far signal imperfectly with residual $y_{\text{res}} = y_2 - (1-\zeta)\sqrt{P_1}h_2\hat{s}_1$ where $\zeta = 0.05$, (3) Decodes own signal from residual with $\gamma_2 = \frac{P_2|h_2|^2}{\zeta P_1|h_2|^2 + \sigma^2}$. Far user decodes directly computing $\gamma_1 = \frac{P_1|h_1|^2}{P_2|h_1|^2 + \sigma^2}$. 

The system calculates rates $R_i = B\log_2(1 + \gamma_i)$ with bandwidth $B = 10$ MHz, checks outage conditions ($\gamma_i < -5$ dB), computes sum-rate $R_{\text{sum}} = R_1 + R_2$, and updates performance metrics.

\subsubsection{DQN Training Workflow}

Figure \ref{fig:dqn_flow} illustrates the complete DQN training procedure.

\begin{figure}[htbp]
    \centering
    \includegraphics[width=0.95\textwidth]{dqn_training.png}
    \caption{Deep Q-Network Training Workflow}
    \label{fig:dqn_flow}
\end{figure}

DQN training initializes Q-network and target network with Dueling architecture, allocates replay buffer for 50,000 transitions, and sets exploration $\epsilon = 1.0$. Each episode starts with environment reset. The training loop observes the 14-dimensional state $s_t$, selects a binary action $a_t \in \{0, 1\}$ using $\epsilon$-greedy exploration (random if $u < \epsilon$, else $\arg\max_a Q(s_t,a;\theta)$). If $a_t = 1$, the file is cached (evicting the least valuable item based on a LFU/LRU heuristic if full); if $a_t = 0$, caching is bypassed. The environment processes the NOMA transmission, calculates reward, observes next state $s_{t+1}$, stores transition $(s_t, a_t, r_t, s_{t+1})$ in buffer with priority $p_i = (|\delta_i| + 10^{-6})^{0.6}$.

When buffer size exceeds 64, samples mini-batch using prioritized probabilities, computes TD targets $y_i = r_i + 0.99\max_{a'} Q(s_{i+1},a';\theta^-)$ using target network, calculates TD errors $\delta_i = y_i - Q(s_i,a_i;\theta)$, updates Q-network minimizing loss $L(\theta) = \frac{1}{64}\sum_{i=1}^{64}\delta_i^2$ using Adam optimizer with learning rate $10^{-3}$, soft updates target network $\theta^- \leftarrow 0.005\theta + 0.995\theta^-$, updates transition priorities based on new TD errors, and decays exploration $\epsilon \leftarrow \max(0.01, 0.995\epsilon)$. Training continues until episode completion, logging metrics (cache hit ratio, reward, loss) after each episode and saving trained model after all 50 episodes.

\subsection{Key Algorithm Designs}

\textbf{Algorithm 4.1: Hybrid Cache-Aware Power Allocation}
\begin{verbatim}
Input:  Users {u1, u2}, Total power P, Cache states {c1, c2}
Output: Power allocation {P1, P2}
1: Sort by channel gain: |h1|^2 <= |h2|^2
2: If c1 == 1 and c2 == 0 Then
3:    P1, P2 <- Compute bounds with zero interference (CIC)
4: Else If c2 == 1 and c1 == 0 Then
5:    P1, P2 <- Compute bounds with perfect SIC (zeta = 0)
6: Else
7:    P1 <- 0.8 * P, P2 <- 0.2 * P  // Baseline fallback
8: Return {P1, P2}
\end{verbatim}

\textbf{Algorithm 4.2: SIC with Imperfection}
\begin{verbatim}
Input:  Signal y2, Powers {P1,P2}, Channels {h1,h2}, zeta=0.05
Output: SINR {gamma1, gamma2}
1: gamma2->1 <- P1|h2|^2/(P2|h2|^2 + sigma^2)  // Decode far
2: y_res <- y2 - (1-zeta)sqrt(P1)*h2*s_hat1     // Cancel
3: gamma2 <- P2|h2|^2/(zeta*P1|h2|^2 + sigma^2) // Decode own
4: gamma1 <- P1|h1|^2/(P2|h1|^2 + sigma^2)      // Far direct
5: Return {gamma1, gamma2}
\end{verbatim}

\textbf{Algorithm 4.3: DQN Cache Update}
\begin{verbatim}
Input:  State st, Q-network Q(s,a;theta), Cache C, epsilon
Output: Action at, Updated cache
1: u <- Random(0,1)
2: If u < epsilon Then at <- RandomAction()
3: Else at <- argmax_a Q(st, a; theta)
4: If at == 1 and file k not in C Then
5:    If |C| == 200 Then
6:       victim <- argmin_{f in C} RequestFrequency(f)
7:       C.remove(victim)
8:    C.add(k)
9: Return at, C
\end{verbatim}

\subsection{Design Considerations}

\textbf{Modularity:} Each component (channel, NOMA, cache, DQN) encapsulates functionality with well-defined interfaces enabling independent testing. NOMA simulator provides \texttt{transmit()} method without exposing SIC internals.

\textbf{Configurability:} All parameters externalize in YAML files (\texttt{config.yaml}) including cache size, Zipf alpha, SIC imperfection, DQN hyperparameters, avoiding hard-coded values.

\textbf{Extensibility:} New cache policies extend \texttt{CachePolicy} base class. Additional channel models inherit from \texttt{ChannelModel} interface. Design supports multi-user NOMA beyond two users.

\textbf{Performance:} NumPy vectorization optimizes critical paths including channel generation and SINR calculations. Multiprocessing enables parallel episode execution for parameter sweeps.

\textbf{Testability:} Unit tests validate modules independently. Mock objects (\texttt{MockChannel}, \texttt{MockCache}) simulate dependencies during testing ensuring correctness.

\thispagestyle{plain}

%-------------------------------------------------------------------------------
%   CHAPTER-V DIVIDER PAGE
%   Matches IOT.pdf format
%   Page numbering continues from previous chapters
%-------------------------------------------------------------------------------

\clearpage

\begin{center}
    \vspace*{10cm}
    
    {\fontsize{16}{19.2}\selectfont\bfseries
    Chapter -- V}
    
    \vspace{1cm}
    
    {\fontsize{16}{19.2}\selectfont\bfseries
    SOFTWARE AND HARDWARE REQUIREMENTS}
    
    \vspace*{\fill}
\end{center}

\thispagestyle{plain}

%-------------------------------------------------------------------------------
%   CHAPTER 5: SOFTWARE AND HARDWARE REQUIREMENTS
%   Following IOT.pdf structure (pages 24-26)
%-------------------------------------------------------------------------------

\clearpage

\section{Software and Hardware Requirements}

\subsection{Software Requirements}

\subsubsection{Software}

\begin{itemize}
    \item Jupyter Notebook / JupyterLab
    \item Visual Studio Code
    \item Google Colab (for cloud-based training)
\end{itemize}

\subsubsection{Operating System}

\begin{itemize}
    \item Windows Operating System (Windows 10/11)
    \item Linux (Ubuntu 20.04 LTS or later - Recommended)
    \item macOS 11 or later
\end{itemize}

\subsubsection{Tools}

\begin{itemize}
    \item Web Browser (Chrome, Firefox, Edge)
    \item Git for version control
    \item GitHub for repository hosting
\end{itemize}

\subsubsection{Core Libraries Used}

NumPy, Pandas, Matplotlib, Seaborn, PyTorch, SciPy, PyYAML, Tqdm

\subsubsection{Detailed Library Description}

\textbf{NumPy}

NumPy is a general-purpose array-processing package. It provides a high-performance multidimensional array object, and tools for working with these arrays. In this project, NumPy is extensively used for channel modeling computations, SINR calculations, power allocation, matrix operations for neural networks, and generating statistical distributions (Rayleigh fading, log-normal shadowing).

\textbf{Pandas}

Pandas is an open-source library that is built on top of NumPy library. It is a Python package that provides various data structures and operations for manipulating numerical data and time series. It is fast and it has high-performance productivity for users. In this project, Pandas is used for storing simulation results in DataFrames, organizing performance metrics (cache hit ratio, outage probability, sum-rate), exporting results to CSV files for analysis, and managing episodic training data from DQN agent.

\textbf{SciPy}

SciPy is a scientific computation library that uses NumPy underneath. It provides more utility functions for optimization, statistical functions, and signal processing. In this project, SciPy is used for generating Zipf distribution for content popularity modeling (alpha = 1.0), Rayleigh distribution for wireless channel fading, statistical analysis and hypothesis testing of simulation results, and signal processing operations for NOMA transmission.

\textbf{Matplotlib}

Matplotlib is a very powerful plotting library useful for those working with Python and NumPy. And for creating statistical inference, it becomes very necessary to visualize our data and Matplotlib is the tool which will be very helpful for this purpose. In this project, Matplotlib is used for plotting training curves (loss, reward, epsilon decay), visualizing cache hit ratio comparison across policies (LRU, LFU, DQN), generating outage probability and sum-rate graphs, creating SINR distribution plots, and producing publication-quality figures for report documentation.

\textbf{PyTorch}

PyTorch is an open-source machine learning library developed by Facebook's AI Research lab. It provides tensor computation with strong GPU acceleration support and deep neural networks built on a tape-based automatic differentiation system. In this project, PyTorch is used for implementing Dueling DQN architecture with separate value and advantage streams, building neural network layers (fully connected, activation functions), automatic differentiation for backpropagation during Q-network training, GPU acceleration for faster training with CUDA support, and implementing experience replay buffer with prioritized sampling.


\subsection{Hardware Requirements}

\subsubsection{Minimum Hardware Specifications}

For basic simulation runs with limited episodes (10-20 episodes):

\begin{table}[htbp]
\centering
\caption{Minimum Hardware Requirements}
\label{tab:min_hardware}
\begin{tabular}{|l|l|}
\hline
\textbf{Component} & \textbf{Specification} \\
\hline
Processor & Intel Core i5 8th gen or AMD Ryzen 5 (4 cores) \\
\hline
RAM & 8 GB DDR4 \\
\hline
Storage & 20 GB available SSD space \\
\hline
Graphics & Integrated graphics (Intel UHD 620 or equivalent) \\
\hline
Network & Internet connection for package installation \\
\hline
\end{tabular}
\end{table}

\subsubsection{Recommended Hardware Specifications}

For comprehensive simulations with large parameter sweeps (50+ episodes, multiple runs):

\begin{table}[htbp]
\centering
\caption{Recommended Hardware Configuration}
\label{tab:rec_hardware}
\begin{tabular}{|l|l|}
\hline
\textbf{Component} & \textbf{Specification} \\
\hline
Processor & Intel Core i7/i9 10th gen or AMD Ryzen 7/9 (8+ cores) \\
\hline
RAM & 16-32 GB DDR4 \\
\hline
Storage & 50 GB NVMe SSD for faster I/O operations \\
\hline
Graphics & NVIDIA GPU (GTX 1660 or better) with CUDA support \\
& for DQN training acceleration \\
\hline
GPU Memory & 4-8 GB VRAM for neural network operations \\
\hline
Network & 10+ Mbps bandwidth for package downloads \\
\hline
\end{tabular}
\end{table}

\subsubsection{GPU Requirements (Optional but Recommended)}

\textbf{NVIDIA GPU with CUDA Support}

For accelerated deep learning training, NVIDIA GPUs with CUDA support significantly reduce training time for DQN models. The Compute Unified Device Architecture (CUDA) is a parallel computing platform and programming model developed by NVIDIA for general computing on GPUs.

Benefits of GPU acceleration for this project:

\begin{itemize}
    \item Faster neural network forward and backward passes during DQN training
    \item Reduced training time from hours to minutes for 50 episodes
    \item Ability to run larger batch sizes (128-256) for more stable training
    \item Parallel processing of multiple simulation episodes
\end{itemize}

Recommended GPU specifications:

\begin{itemize}
    \item NVIDIA GPU with compute capability 3.5 or higher
    \item CUDA Toolkit 11.3 or later
    \item cuDNN 8.2 or later for optimized neural network operations
    \item Minimum 4 GB VRAM, 8+ GB recommended for large-scale experiments
\end{itemize}


\thispagestyle{plain}

%-------------------------------------------------------------------------------
%   CHAPTER-VI DIVIDER PAGE
%   Matches IOT.pdf format
%   Page numbering continues from previous chapters
%-------------------------------------------------------------------------------

\clearpage

\begin{center}
    \vspace*{10cm}
    
    {\fontsize{16}{19.2}\selectfont\bfseries
    Chapter -- VI}
    
    \vspace{1cm}
    
    {\fontsize{16}{19.2}\selectfont\bfseries
    CODE TEMPLATES AND IMPLEMENTATION DETAILS}
    
    \vspace*{\fill}
\end{center}

\thispagestyle{plain}

%-------------------------------------------------------------------------------
%   CHAPTER 6: CODE TEMPLATES AND IMPLEMENTATION DETAILS
%   Following IOT.pdf pages 27-40 structure
%-------------------------------------------------------------------------------

\clearpage

\section{Code Templates and Implementation Details}

This chapter presents detailed implementation of key modules with comprehensive documentation of classes, methods, input parameters, and output specifications corresponding to the updated Cache-Aided NOMA framework.

\subsection{Channel and NOMA Modeling Module}

\subsubsection{Class: TimeVaryingChannel}

\textbf{Purpose:} Implements mixed fading channels (Rayleigh/Rician) and computes realistic time-varying channel gains incorporating path loss and shadowing.

\textbf{Attributes:}
\begin{itemize}
    \item \texttt{n\_users}: Number of users in the cell.
    \item \texttt{r\_max}: Maximum cell radius in meters.
\end{itemize}

\textbf{Methods:}
\begin{verbatim}
class TimeVaryingChannel:
    def __init__(self, carrier_freq=28e9, velocity=3.0, sampling_interval=0.001):
        '''
        Input:
            carrier_freq (float): Carrier frequency (Hz)
            velocity (float): User velocity (m/s)
            sampling_interval (float): Sampling interval (s)
        '''
        self.carrier_freq = carrier_freq
        self.velocity = velocity
        self.doppler_freq = velocity * carrier_freq / 3e8
    
    def generate_time_series(self, num_samples, fading_type='rayleigh'):
        '''
        Computes composite time-correlated fading samples 
        using Jake's model.
        
        Output:
            ndarray: Time-correlated fading gains
        '''
        # Implementation of Jake's model AR(1) process...
\end{verbatim}

\subsubsection{Module: power\_allocation.py}

\textbf{Purpose:} Implements the Hybrid SIC/CIC power allocation logic dynamically adapting to the state of edge caches.

\textbf{Methods:}
\begin{verbatim}
def allocate_power_cache_aware(gain_w, gain_s, cfg, 
                               weak_cached, strong_cached):
    '''
    Allocate power intelligently based on cache availability 
    (enabling hybrid CIC).
    
    Input:
        gain_w, gain_s (float): Channel gains for weak/strong
        cfg (Config): System configuration with power bounds
        weak_cached, strong_cached (bool): Cache indicators
        
    Output:
        p_w, p_s (float): Allocated transmit powers
        ok (bool): Feasibility indicator
        info (dict): Telemetry metadata (e.g., CIC usage)
    '''
    if weak_cached and not strong_cached:
        # Scale back weak power (CIC)
    elif strong_cached and not weak_cached:
        # Perfect SIC condition (zeta=0)
    # ... calculates bounds and returns optimized power pair
    return p_w, p_s, ok, info
\end{verbatim}

\subsection{Caching Module}

\subsubsection{Class: DQNCache}

\textbf{Purpose:} Implements the Double Dueling DQN agent with a binary action space and a CIC-aware reward mechanism for proactive cache placement.

\textbf{Attributes:}
\begin{itemize}
    \item \texttt{capacity}: Maximum number of files in the cache.
    \item \texttt{state\_dim}: 14-dimensional NOMA-aware state.
    \item \texttt{action\_dim}: 2 (Binary Action Space).
\end{itemize}

\textbf{Methods:}
\begin{verbatim}
class DQNCache(CacheBase):
    def __init__(self, capacity: int, num_files: int, 
                 num_users: int, learning_rate: float = 1e-3, 
                 gamma: float = 0.99):
        '''
        Initializes the Double Dueling DQN network, prioritized 
        replay buffer, and 14-dimensional state tracking.
        '''
        super().__init__(capacity)
        self.state_dim = 14
        self.action_dim = 2 # Binary: 0 (Reject), 1 (Cache)
        self.memory = PrioritizedReplayBuffer(maxlen=50000)

        self._build_model()
        
    def _build_model(self):
        '''
        Builds the Dueling architecture dividing into Value 
        and Advantage streams.
        '''
        pass
        
    def act(self, state):
        '''
        Epsilon-greedy selection of binary action.
        
        Input: state (ndarray) 14-dim state vector.
        Output: action (int) 0 or 1.
        '''
        if random.random() < self.epsilon:
            return random.randint(0, 1)
        state_tensor = torch.FloatTensor(state)
        q_values = self.q_network(state_tensor)
        return q_values.argmax().item()
\end{verbatim}

\thispagestyle{plain}

%-------------------------------------------------------------------------------
%   CHAPTER-VII DIVIDER PAGE
%-------------------------------------------------------------------------------

\clearpage

\begin{center}
    \vspace*{10cm}
    
    {\fontsize{16}{19.2}\selectfont\bfseries
    Chapter -- VII}
    
    \vspace{1cm}
    
    {\fontsize{16}{19.2}\selectfont\bfseries
    EXPERIMENTAL RESULTS AND ANALYSIS}
    
    \vspace*{\fill}
\end{center}

\thispagestyle{plain}

%-------------------------------------------------------------------------------
%   CHAPTER 7: EXPERIMENTAL RESULTS AND ANALYSIS
%-------------------------------------------------------------------------------

\clearpage

\section{Experimental Results and Analysis}

This chapter details the performance evaluation of the proposed Cache-Aided NOMA (CA-NOMA with CIC) framework against conventional baselines. The analysis is divided into fundamental performance bounds, Deep Reinforcement Learning (DRL) cache efficiency, training convergence stability, and a comprehensive Monte Carlo policy verification across 9 critical network metrics.

\subsection{Baseline Comparison: CE-NOMA vs CA-NOMA with CIC}

We first evaluate the fundamental throughput and reliability improvements when caching is introduced to the NOMA architecture. Here, \textbf{CE-NOMA} (Cache-Enabled NOMA) represents a standard NOMA system equipped with a local cache storing popular files, without interference cancellation benefits. \textbf{CA-NOMA with CIC} (Cache-Aided NOMA) represents our proposed architecture where the cache directly assists the physical layer by enabling Cache-Aided Interference Cancellation (CIC).

\begin{figure}[htbp]
\centerline{\includegraphics[width=0.75\textwidth]{results/base_comparison/fig_sum_rate_vs_snr.png}}
\caption{System Sum-Rate vs. SNR (dB) comparing Conventional NOMA, CE-NOMA, and CA-NOMA with CIC.}
\label{fig:base_sum_rate}
\end{figure}

\textbf{Sum-Rate Expansion:} As illustrated in Fig. \ref{fig:base_sum_rate}, the integration of proactive caching immediately shifts the sum-rate ceiling. While conventional NOMA saturates rapidly due to heavy co-channel interference, CE-NOMA yields moderate improvements purely via traffic offloading. However, the CA-NOMA with CIC framework dynamically exploits cache hits to eliminate successive interference cancellation (SIC) error floors ($\zeta$), driving the sum-rate boundary exponentially higher under identical channel conditions.

\begin{figure}[htbp]
\centerline{\includegraphics[width=0.75\textwidth]{results/base_comparison/fig_outage_vs_snr.png}}
\caption{Outage Probability vs. SNR (dB) highlighting extreme reliability gains under CIC.}
\label{fig:base_outage}
\end{figure}

\textbf{Outage Probability Mitigation:} Fig. \ref{fig:base_outage} demonstrates the logarithmic decay of outage probabilities. The CA-NOMA with CIC architecture achieves a strict 1-2 order of magnitude collapse in network outage events, validating that resolving interference mathematically at the edge significantly suppresses bit-error propagation for both strong and weak users simultaneously.

\begin{figure}[htbp]
\centerline{\includegraphics[width=0.75\textwidth]{results/base_comparison/fig_cache_hit_vs_size.png}}
\caption{Cache Hit Ratio vs. Cache Size for Top-K analytical limit vs dynamic tracking.}
\label{fig:base_hit_rate}
\end{figure}

\textbf{Cache Hit Rate:} Fig. \ref{fig:base_hit_rate} outlines the fundamental cache hit ratio scaling against physical storage capacity ($C$). Under a Zipf distribution ($\alpha=1.0$), localized storage rapidly absorbs up to 70\% of network requests with only 10\% catalog capacity (200 files), proving the feasibility of edge-deployed storage modules.

\subsection{DRL Caching Performance (DQN Comparison)}

While static policies (like Top-K) assume perfect prior knowledge of traffic, our proposed system employs a Double Dueling Deep Q-Network (DQN) to blindly learn and adapt to request patterns.

\begin{figure}[htbp]
\centerline{\includegraphics[width=0.75\textwidth]{results/dqn_comparison/fig_sum_rate_vs_snr.png}}
\caption{DQN-driven Sum-Rate vs. SNR compared to traditional LFU/LRU heuristics.}
\label{fig:dqn_sum_rate}
\end{figure}

\textbf{DQN Sum-Rate Optimization:} Fig. \ref{fig:dqn_sum_rate} confirms that the DQN agent successfully matches the theoretical maximums established by the omniscient Top-K oracle. By explicitly encoding a CIC-aware reward mechanism, the agent actively prefers caching payloads that maximize multi-user interference cancellation, yielding higher sum-rates than both LRU (Least Recently Used) and LFU (Least Frequently Used) heuristics.

\begin{figure}[htbp]
\centerline{\includegraphics[width=0.75\textwidth]{results/dqn_comparison/fig_outage_vs_snr.png}}
\caption{DQN Outage Probability showcasing stability against random eviction policies.}
\label{fig:dqn_outage}
\end{figure}

\begin{figure}[htbp]
\centerline{\includegraphics[width=0.75\textwidth]{results/dqn_comparison/fig_cache_hit_vs_size.png}}
\caption{DQN Hit Rate evaluated at $C=200$, bridging the gap between LFU and Top-K.}
\label{fig:dqn_hit_rate}
\end{figure}

\textbf{DQN Hit Rate Autonomy:} As shown in Fig. \ref{fig:dqn_hit_rate}, evaluated specifically at the hardware constraint $C=200$, the DQN framework achieves an absolute hit rate of $\approx 70.38\%$. This trails the impossible mathematical Top-K threshold (70.92\%) by a negligible margin, while substantially outperforming temporal-biased LRU policies.

\subsection{Deep Q-Network Training and Evaluation}

To guarantee real-world deployment feasibility, the DQN architecture's stability must be verified across extensive episodes.

\begin{figure}[htbp]
\centerline{\includegraphics[width=0.85\textwidth]{results/dqn_training_curves.png}}
\caption{DQN Training Curves over 500 episodes displaying Hit Rate convergence, Exploration ($\epsilon$) decay, and Huber Loss minimization.}
\label{fig:training_curves}
\end{figure}

\textbf{Convergence Stability:} Fig. \ref{fig:training_curves} traces the Dueling DQN's training topology. The cache hit rate strictly plateaus near episode 150, corresponding identically with the decay of the exploration factor ($\epsilon$). The Huber Loss maps the characteristic double-spike pattern of Prioritized Experience Replay (PER), proving that the agent rapidly isolates high-value state-action pairs without suffering from catastrophic forgetting.

\subsection{Comprehensive Policy Verification (Monte Carlo Analysis)}

A rigorous Monte Carlo simulation mapping 1,000 statistical realizations was executed to isolate the performance bounds of CA-NOMA with CIC across all evaluated caching methodologies. The large 9-subplot comparative matrix generated by the simulation framework has been broken down into independent figures for detailed analysis.

\begin{figure}[htbp]
\centerline{\includegraphics[width=0.55\textwidth]{results/split_plots/fig_sum_rate.png}}
\caption{Monte Carlo: Sum-Rate Distribution}
\label{fig:split_sum_rate}
\end{figure}

\textbf{1. Sum-Rate Distribution:} The comprehensive variance evaluation (Fig. \ref{fig:split_sum_rate}) confirms the DQN protocol consistently enforces a massive throughput explosion over No-Cache baselines, tightly trailing the theoretical Top-K oracle.

\begin{figure}[htbp]
\centerline{\includegraphics[width=0.55\textwidth]{results/split_plots/fig_outage.png}}
\caption{Monte Carlo: Outage Probability}
\label{fig:split_outage}
\end{figure}

\textbf{2. Outage Probability:} As seen in Fig. \ref{fig:split_outage}, outage variances heavily compress under DQN influence, ensuring strict transmission reliability specifically within highly volatile -10 to 5 dB SNR regions.

\begin{figure}[htbp]
\centerline{\includegraphics[width=0.55\textwidth]{results/split_plots/fig_cache_hit.png}}
\caption{Monte Carlo: Cache Hit Ratio}
\label{fig:split_hit}
\end{figure}

\textbf{3. Cache Hit Ratio Consistency:} The tight confidence intervals surrounding the DQN hit rate (70.38\%) in Fig. \ref{fig:split_hit} prove the deep reinforcement agent accurately models the Zipf distribution structure dynamically, tracking popularity shifts autonomously.

\begin{figure}[htbp]
\centerline{\includegraphics[width=0.55\textwidth]{results/split_plots/fig_weak_rate.png}}
\caption{Monte Carlo: Weak User Rate}
\label{fig:split_weak}
\end{figure}

\textbf{4. Weak User Delivery:} Weak user bounds (Fig. \ref{fig:split_weak}) specifically benefit from CA-NOMA. By retrieving weak-user payloads from the cache, the base station commits 100\% of physical wireless power to the strong user, inadvertently multiplying weak user effective rates via interference elimination.

\begin{figure}[htbp]
\centerline{\includegraphics[width=0.55\textwidth]{results/split_plots/fig_strong_rate.png}}
\caption{Monte Carlo: Strong User Rate}
\label{fig:split_strong}
\end{figure}

\textbf{5. Strong User Delivery:} As mapped in Fig. \ref{fig:split_strong}, the strong user, previously handicapped by imperfect SIC residuals ($\zeta=0.05$), achieves flawless decoding when the interfering weak payload is localized in cache memory, resulting in steep linear rate scaling.

\begin{figure}[htbp]
\centerline{\includegraphics[width=0.55\textwidth]{results/split_plots/fig_fairness.png}}
\caption{Monte Carlo: Jain's Fairness Index}
\label{fig:split_fairness}
\end{figure}

\textbf{6. System Fairness:} Jain's Fairness Index (Fig. \ref{fig:split_fairness}) demonstrates that CA-NOMA perfectly equalizes resource distribution, yielding fairness indices near 0.95, far superseding traditional power-domain NOMA inequalities.

\begin{figure}[htbp]
\centerline{\includegraphics[width=0.55\textwidth]{results/split_plots/fig_cic_rate.png}}
\caption{Monte Carlo: CIC Exploitation Rate}
\label{fig:split_cic}
\end{figure}

\textbf{7. CIC Efficiency Dominance:} Evaluating NOMA interference boundaries in Fig. \ref{fig:split_cic}, the DQN framework maps an average CIC exploitation rate of 3.776, an irrefutable 5.28x improvement over static Top-K, solidifying that the agent successfully learns to cache files that specifically trigger mathematical cancellation.

\begin{figure}[htbp]
\centerline{\includegraphics[width=0.55\textwidth]{results/split_plots/fig_energy_eff.png}}
\caption{Monte Carlo: Energy Efficiency}
\label{fig:split_energy}
\end{figure}

\textbf{8. Energy Efficiency:} As shown in Fig. \ref{fig:split_energy}, localized cached payload deliveries bypass base-station retransmission cycling, spiking energy boundaries to a peak evaluation of 18.66 bits/Joule natively surpassing dynamic LRU routes.

\begin{figure}[htbp]
\centerline{\includegraphics[width=0.55\textwidth]{results/split_plots/fig_throughput.png}}
\caption{Monte Carlo: Average Throughput}
\label{fig:split_throughput}
\end{figure}

\textbf{9. Average Throughput:} The composite throughput geometry (Fig. \ref{fig:split_throughput}) definitively confirms that CA-NOMA with CIC under Deep RL supervision fundamentally bridges the gap between hardware interference limits and perfect channel capacities, outperforming non-cached models geometrically.

%=============================================================
% VII. CONCLUSION
%=============================================================

\section{Conclusion}

This paper has definitively formalized a Double Dueling Deep Q-Network pipeline structured specifically around proactively generating localized caching policies mapped mathematically into NOMA network transmission paths for dynamic 6G infrastructure. Replacing infinite algorithmic decision boundaries with a strict binary resolution model interacting sequentially with a heavily condensed 14-dimensional channel state, the module autonomously isolates optimal delivery sequences prioritizing the exclusive reconstruction elements driving Cache-Aided Interference Cancellation (CIC). The proposed hybrid SIC/CIC framework enables both weak-user and strong-user interference cancellation modes, while the cache-aware dynamic power allocation strategy adapts transmit power coefficients in real-time to exploit the altered interference structure---delivering additional energy savings and throughput gains beyond what static power allocation can achieve. Independent Monte Carlo validations effectively verify sustained performance evaluations matching explicit ideal oracle limitations down to 0.54\% hit probability variances. Simultaneously isolating multi-user interplay rewards results directly inside unprecedented 5.28x factor CIC expansions. These properties effectively guarantee comprehensive reductions dictating absolute operational downtime (reducing outage thresholds uniformly 14x against base limitations), firmly justifying Deep RL frameworks acting centrally inside interconnected 6G edge communication layers.

Future exploration trajectories anticipate formulating scaled reinforcement interactions targeting dynamic user mobility topologies combined intimately via localized federated network coordination deployments spanning wide-area multi-cell coverage bounds.


%=============================================================
% ACKNOWLEDGMENT
%=============================================================
\section*{Acknowledgment}
The authors respectfully acknowledge the critical project steering and invaluable supervisory oversight uniquely contributed by Dr. Biswajit Ghosh representing the Department of Information Technology situated within Techno Main Salt Lake, directly facilitating the core developments powering this Project.


%=============================================================
% REFERENCES
%=============================================================

\section*{References}
\addcontentsline{toc}{section}{References}
\justifying

\begin{thebibliography}{99}

\bibitem{b_breslau_zipf}
L. Breslau, P. Cao, L. Fan, G. Phillips, and S. Shenker, ``Web caching and Zipf-like distributions: Evidence and implications,'' in \textit{Proc. IEEE INFOCOM}, vol. 1, pp. 126--134, 1999.

\bibitem{b_golrezaei}
N. Golrezaei, A. F. Molisch, A. G. Dimakis, and G. Caire, ``Femtocaching and device-to-device collaboration: A new architecture for wireless video distribution,'' \textit{IEEE Commun. Mag.}, vol. 51, no. 4, pp. 142--149, Apr. 2013.

\bibitem{b_maddah_ali}
M. A. Maddah-Ali and U. Niesen, ``Fundamental limits of caching,'' \textit{IEEE Trans. Inf. Theory}, vol. 60, no. 5, pp. 2856--2867, May 2014.

\bibitem{b_dai_noma}
L. Dai, B. Wang, Y. Yuan, S. Han, C.-L. I, and Z. Wang, ``Non-orthogonal multiple access for 5G: Solutions, challenges, opportunities, and future research trends,'' \textit{IEEE Commun. Mag.}, vol. 53, no. 9, pp. 74--81, Sep. 2015.

\bibitem{b_ji_d2d}
M. Ji, G. Caire, and A. F. Molisch, ``Wireless device-to-device caching networks: Basic principles and system performance,'' \textit{IEEE J. Sel. Areas Commun.}, vol. 34, no. 1, pp. 176--189, Jan. 2016.

\bibitem{b_guo_edge}
S. Guo, B. Xiao, Y. Yang, and Y. Yang, ``Energy-efficient dynamic offloading and resource scheduling in mobile cloud computing,'' in \textit{Proc. IEEE INFOCOM}, pp. 1--10, 2016.

\bibitem{b_ding_survey}
Z. Ding, X. Lei, G. K. Karagiannidis, R. Schober, J. Yuan, and V. K. Bhargava, ``A survey on non-orthogonal multiple access for 5G networks: Research challenges and future trends,'' \textit{IEEE J. Sel. Areas Commun.}, vol. 35, no. 10, pp. 2181--2195, Oct. 2017.

\bibitem{b_zhang_stochastic}
H. Zhang, Y. Qiu, X. Chu, K. Long, and V. C. M. Leung, ``Fog radio access networks: Mobility management, interference mitigation, and resource optimization,'' \textit{IEEE Wireless Commun.}, vol. 24, no. 6, pp. 120--127, Dec. 2017.

\bibitem{b_liu_uav}
Y. Liu, Z. Qin, M. Elkashlan, Z. Ding, A. Nallanathan, and L. Hanzo, ``Nonorthogonal multiple access for 5G and beyond,'' \textit{Proc. IEEE}, vol. 105, no. 12, pp. 2347--2381, Dec. 2017.

\bibitem{b_li_drl}
R. Li, C. Zhao, S. Guo, T. Zhang, and M. Liu, ``Deep reinforcement learning for resource management in network slicing,'' \textit{IEEE Access}, vol. 6, pp. 74429--74441, 2018.

\bibitem{b_cao_joint}
B. Cao, L. Zhang, Y. Li, D. Feng, and W. Cao, ``Intelligent offloading in multi-access edge computing: A state-of-the-art review and framework,'' \textit{IEEE Commun. Mag.}, vol. 57, no. 3, pp. 56--62, Mar. 2019.

\bibitem{b_elghariani_pred}
A. Elghariani and D. Zeghlache, ``A learning-based adaptive caching framework for mobile content delivery,'' in \textit{Proc. IEEE GLOBECOM}, pp. 1--6, 2019.

\bibitem{b_chen_ml}
M. Chen, W. Saad, C. Yin, and M. Debbah, ``Data correlation-aware resource management in wireless powered federated edge learning,'' \textit{IEEE Trans. Wireless Commun.}, vol. 20, no. 3, pp. 1498--1512, Mar. 2021.

\bibitem{b_survey_bepari}
S. Bepari et al., ``A survey on cache-aided NOMA,'' \textit{IEEE Communications Surveys \& Tutorials}, 2023.

\bibitem{b_bepari_partial}
S. Bepari et al., ``Partial cache-aided NOMA,'' \textit{IEEE Transactions on Wireless Communications}, 2024.

\bibitem{b_wang_dueling}
Z. Wang, T. Schaul, M. Hessel, H. van Hasselt, M. Lanctot, and N. de Freitas, ``Dueling network architectures for deep reinforcement learning,'' in \textit{Proc. ICML}, pp. 1995--2003, 2016.

\bibitem{b_van_hasselt}
H. van Hasselt, A. Guez, and D. Silver, ``Deep reinforcement learning with double Q-learning,'' in \textit{Proc. AAAI}, pp. 2094--2100, 2016.

\bibitem{b_schaul_per}
T. Schaul, J. Quan, I. Antonoglou, and D. Silver, ``Prioritized experience replay,'' in \textit{Proc. ICLR}, 2016.

\end{thebibliography}

\thispagestyle{plain}

%-------------------------------------------------------------------------------
%   APPENDICES
%-------------------------------------------------------------------------------

\clearpage

\section*{Appendix A: Mathematical Derivations}
\addcontentsline{toc}{section}{Appendix A: Mathematical Derivations}

\subsection*{A.1 SINR Derivation for NOMA Users}

Consider two users paired for NOMA transmission: far user $U_F$ (weak channel) and near user $U_N$ (strong channel) with channel coefficients $h_F$ and $h_N$ respectively, where $|h_F|^2 < |h_N|^2$. The base station transmits the superposed signal:
\begin{equation}
x = \sqrt{P_F} s_F + \sqrt{P_N} s_N
\end{equation}
where $P_F = 0.8P_{\text{total}}$ and $P_N = 0.2P_{\text{total}}$ are power allocations following NOMA principle (more power to weaker user), and $s_F$, $s_N$ are unit-power symbols.

\textbf{Near User (Strong Channel):} Receives signal:
\begin{equation}
y_N = h_N x + n_N = h_N(\sqrt{P_F} s_F + \sqrt{P_N} s_N) + n_N
\end{equation}

Following SIC protocol, near user first decodes far user's signal (treating own signal as interference):
\begin{equation}
\gamma_{F \rightarrow N} = \frac{P_F |h_N|^2}{P_N |h_N|^2 + N_0}
\end{equation}

After canceling far user's signal with imperfection factor $\epsilon = 0.05$, near user decodes own signal:
\begin{equation}
\gamma_N = \frac{P_N |h_N|^2}{\epsilon P_F |h_N|^2 + N_0}
\end{equation}

\textbf{Far User (Weak Channel):} Directly decodes own signal, treating near user as noise:
\begin{equation}
\gamma_F = \frac{P_F |h_F|^2}{P_N |h_F|^2 + N_0}
\end{equation}

Outage occurs when $\gamma < \gamma_{\text{target}}$, where $\gamma_{\text{target}} = 2^{R_{\text{target}}} - 1$ for target rate $R_{\text{target}} = 0.5$ bps/Hz.

\subsection*{A.2 Cache Hit Ratio Under Zipf Distribution}

For content popularity following Zipf distribution with skewness parameter $\alpha = 1.0$, the probability of requesting the $i$-th most popular file is:
\begin{equation}
p_i = \frac{i^{-\alpha}}{\sum_{j=1}^{N} j^{-\alpha}} = \frac{1/i}{H_N}
\end{equation}
where $N = 2000$ is the catalog size and $H_N = \sum_{j=1}^{N} 1/j \approx \ln(N) + \gamma$ is the harmonic number ($\gamma \approx 0.577$ is Euler's constant).

For cache capacity $K = 200$ storing top-$K$ files, theoretical hit ratio:
\begin{equation}
H_{\text{theory}} = \sum_{i=1}^{K} p_i = \frac{H_K}{H_N} = \frac{\sum_{i=1}^{200} 1/i}{\sum_{j=1}^{2000} 1/j} \approx \frac{5.878}{8.178} \approx 0.719
\end{equation}

For $\alpha = 1.0$ and $K/N = 0.1$, theoretical hit ratio is approximately 71.9\%. Practical implementations achieve lower values (38-44\%) due to temporal variations and request patterns.

\subsection*{A.3 DQN Bellman Optimality Equation}

The DQN agent learns optimal cache placement policy by approximating the action-value function $Q(s, a)$. The Bellman optimality equation:
\begin{equation}
Q^*(s, a) = \mathbb{E}_{s'}\left[r(s, a) + \gamma \max_{a'} Q^*(s', a')\right]
\end{equation}
where $s$ is state (cache contents, request history), $a$ is action (content to cache), $r(s,a)$ is immediate reward, $\gamma = 0.95$ is discount factor, and $s'$ is next state.

The neural network minimizes the loss:
\begin{equation}
L(\theta) = \mathbb{E}_{(s,a,r,s') \sim \mathcal{D}}\left[\left(r + \gamma \max_{a'} Q(s', a'; \theta^-) - Q(s, a; \theta)\right)^2\right]
\end{equation}
where $\theta$ are network parameters, $\theta^-$ are target network parameters updated via soft update: $\theta^- \leftarrow \tau\theta + (1-\tau)\theta^-$ with $\tau = 0.005$.

\clearpage

\section*{Appendix B: Configuration Files}
\addcontentsline{toc}{section}{Appendix B: Configuration Files}

\subsection*{B.1 Main Simulation Configuration (config.py)}

\begin{verbatim}
# Simulation Configuration for Cache-Aided NOMA

# Random seed
RANDOM_SEED = 2025

# Content catalog
NUM_FILES = 2000              # Total unique files
ZIPF_ALPHA = 1.0              # Zipf skewness parameter

# Users and requests
NUM_USERS = 200               # Users in cell
REQUESTS_PER_USER = 50        # Requests per user

# Cache configuration
CACHE_SIZE = 200              # Cache capacity (10% of catalog)
CACHE_POLICY = "stable_dqn"   # Policy: topk/lru/lfu/stable_dqn
CACHE_DELIVERY_RATE = 10.0    # bits/s/Hz

# NOMA parameters
TX_POWER = 1.0                # Total transmit power (W)
NOISE_POWER = 1e-9            # Noise power (W)
CELL_RADIUS = 500.0           # Cell radius (meters)
PATHLOSS_EXPONENT = 3.5       # Path loss exponent
MIN_DISTANCE = 1.0            # Minimum user distance (m)

PAIR_USERS = True             # Enable NOMA pairing
PAIRING_METHOD = "extreme"    # Pairing method
POWER_COEFF_WEAK = 0.8        # Power for far user
POWER_COEFF_STRONG = 0.2      # Power for near user
SIC_IMPERFECTION = 0.05       # Residual interference (5%)
TARGET_RATE_BPS = 0.5         # Target rate (bps/Hz)

# DQN Training Configuration
RL_TRAINING_EPISODES = 50     # Training episodes
RL_STEPS_PER_EPISODE = 1000   # Steps per episode
RL_TRAINING_STEPS = 50000     # Total training steps

# Epsilon-greedy exploration
RL_EPSILON_START = 1.0        # Start with full exploration
RL_EPSILON_END = 0.01         # End with 1% exploration
RL_EPSILON_DECAY_STEPS = 25000  # Decay over first half

# Neural network architecture
RL_HIDDEN_DIMS = [128, 64]    # Hidden layer sizes
RL_LEARNING_RATE = 0.0001     # Learning rate (Adam optimizer)
RL_GAMMA = 0.95               # Discount factor
RL_BATCH_SIZE = 64            # Training batch size
RL_REPLAY_BUFFER_SIZE = 50000 # Experience replay buffer

# Training stability
RL_GRADIENT_CLIP = 10.0       # Gradient clipping max norm
RL_TAU = 0.005                # Soft target update rate
RL_TRAIN_FREQUENCY = 4        # Train every N steps

# Reward function
RL_REWARD_CACHE_HIT = 10.0
RL_REWARD_CACHE_MISS_SUCCESS = -1.0
RL_REWARD_NOMA_FAILURE = -5.0
RL_REWARD_OUTAGE = -5.0

# Prioritized replay
RL_USE_PRIORITIZED_REPLAY = True
RL_PRIORITY_ALPHA = 0.6
RL_PRIORITY_BETA = 0.4

# Evaluation
RL_EVAL_REQUESTS = 5000       # Requests per evaluation
NUM_RUNS = 10                 # Monte Carlo runs
\end{verbatim}

\clearpage

\section*{Appendix C: Experimental Results Data}
\addcontentsline{toc}{section}{Appendix C: Experimental Results Data}

\subsection*{C.1 Performance Metrics Summary}

\begin{verbatim}
===============================================================================
         CACHE-AIDED NOMA EXPERIMENTAL RESULTS
===============================================================================
Date: December 2025
Configuration: Stable DQN with 50 training episodes

--- CACHE PERFORMANCE ---
Policy: Stable DQN
Cache Hit Ratio:         0.442 ± 0.018  (44.2%)
Cache Miss Ratio:        0.558 ± 0.018  (55.8%)
Average Access Time:     42.5 ± 2.8 ms
Cached Content Latency:  2.0 ms
Backhaul Fetch Latency:  100.0 ms

--- NOMA PHYSICAL LAYER PERFORMANCE ---
System Sum Rate:         52.3 ± 3.2 Mbps
Near User Avg Rate:      34.8 Mbps
Far User Avg Rate:       17.5 Mbps
Near User SINR (median): 19.4 dB
Far User SINR (median):  7.8 dB
NOMA Success Rate:       68.5%

--- OUTAGE STATISTICS ---
Overall Outage Probability: 8.2 ± 1.2%
Near User Outage:           2.1%
Far User Outage:            14.3%
Cell-Edge Outage (>450m):   18.7%
Target Rate:                0.5 bps/Hz

--- DQN TRAINING METRICS ---
Training Episodes:       50
Training Steps:          50,000
Training Time:           45.3 minutes
Final Epsilon:           0.010
Cumulative Reward:       +2847.3
Average Loss:            0.00234
Convergence Episode:     38

--- BASELINE COMPARISON ---

Policy          Hit Rate    Outage     Latency    Improvement
----------------------------------------------------------------
Top-K           0.380       0.105      48.2 ms    Baseline
LRU             0.310       0.128      58.6 ms    -18.4%
LFU             0.325       0.120      55.3 ms    -14.5%
Stable DQN      0.442       0.082      42.5 ms    +16.3%

DQN vs Top-K Improvements:
  Hit Rate:       +16.3%
  Outage:         -21.9% (reduction)
  Latency:        -11.8% (reduction)
  NOMA Success:   +37.0%

--- STATISTICAL SIGNIFICANCE ---
Paired t-test (DQN vs Top-K):
  Hit Rate:    p < 0.001  (highly significant)
  Outage:      p < 0.001  (highly significant)
  
--- COMPUTATIONAL PERFORMANCE ---
DQN Inference Time:      2.1 ms per decision
Memory Usage:            ~450 MB (replay buffer)
Model Size:              1.2 MB (saved checkpoint)

--- SIMULATION PARAMETERS ---
Total Requests Evaluated: 50,000 (10 runs × 5000 requests)
Content Catalog:          2000 files
Cache Capacity:           200 files (10%)
User Count:               200 users
Cell Radius:              500 meters
Zipf Alpha:               1.0
===============================================================================
\end{verbatim}

\subsection*{C.2 Per-Episode Training Log (Sample)}

\begin{verbatim}
Episode | Steps | Epsilon | Hit Rate | Reward  | Loss     | Time
--------|-------|---------|----------|---------|----------|-------
1       | 1000  | 0.980   | 0.182    | -234.5  | 0.12340  | 52s
5       | 5000  | 0.920   | 0.215    | -145.2  | 0.08721  | 48s
10      | 10000 | 0.840   | 0.268    | -52.8   | 0.05432  | 51s
15      | 15000 | 0.760   | 0.311    | +23.1   | 0.03214  | 49s
20      | 20000 | 0.680   | 0.358    | +98.4   | 0.01987  | 47s
25      | 25000 | 0.600   | 0.389    | +156.7  | 0.01234  | 50s
30      | 30000 | 0.010   | 0.412    | +198.2  | 0.00876  | 46s
35      | 35000 | 0.010   | 0.431    | +223.5  | 0.00543  | 48s
40      | 40000 | 0.010   | 0.438    | +241.8  | 0.00321  | 47s
45      | 45000 | 0.010   | 0.441    | +252.4  | 0.00267  | 49s
50      | 50000 | 0.010   | 0.443    | +258.1  | 0.00234  | 48s

Note: Epsilon decays linearly from 1.0 to 0.01 over episodes 1-25,
      then remains at 0.01 for episodes 26-50 (exploitation phase).
\end{verbatim}

\thispagestyle{plain}

\end{document}

